%!TeX program = xelatex
\documentclass[12pt,hyperref,a4paper,UTF8]{ctexart}
\usepackage{SJTUReport}
\usepackage{graphicx}
\usepackage{booktabs}
\usepackage{float}
\usepackage{longtable}
\usepackage{array}
\usepackage{amsmath}

%%-------------------------------正文开始---------------------------%%
\begin{document}

%%-----------------------封面--------------------%%
%% 封面信息（如标题、姓名、学号等）请在 SJTUReport.sty 文件中修改
\cover
\clearpage
%%------------------摘要-------------%%
\begin{abstract}
本文围绕电商客户行为分析系统中的实时处理与流式分析模块展开研究与实现。针对传统批处理分析只能反映静态历史数据、难以及时捕捉最新交易动态的问题，本文独立设计并实现了一套基于 Spark Structured Streaming 的完整实时数据处理链路。该链路采用模块化架构，由事件生成器（\texttt{realtime\_producer.py}）、流处理引擎（\texttt{realtime\_processor.py}）和报表生成器（\texttt{realtime\_reporting.py}）三个核心组件协同构成，覆盖了从模拟交易事件生成、流式数据接入、微批次处理、多维聚合统计、风险客户识别到可视化仪表盘输出的全流程。

在事件生成层，本文设计了基于种子数据受控变异的模拟方法，以原始电商客户行为数据集为种子池，对消费金额、购买件数、评分、距上次购买天数等关键字段施加受控随机扰动，并动态推导满意度等级，从而兼顾模拟事件的统计真实性与业务逻辑自洽性。在流处理层，本文基于 Spark Structured Streaming 的 \texttt{foreachBatch} 机制构建了微批次分析引擎，在每个微批次中依次完成事件留存、城市维度聚合、会员维度聚合、风险客户识别和 KPI 指标汇总，并设计了增强型风险客户判定规则——综合距上次购买天数、满意度等级和平均评分三维度进行判定，相比单条件规则具有更高的检测灵敏度。在结果展示层，本文设计并实现了无外部依赖的 HTML 实时仪表盘和 Markdown 摘要，将结构化 CSV 分析结果转化为直观的可视化报告。

系统同时支持文件流模式和 Kafka 模式两种运行方式，通过统一 Schema 设计实现两种模式的处理逻辑复用。文件流模式降低了教学演示场景下的环境依赖与部署复杂度；Kafka 模式则预留了向生产级消息队列式实时处理架构扩展的完整能力。Docker Compose 通过 profile 机制按需编排 Kafka 与 Zookeeper 服务，兼顾了默认运行的简洁性与高阶场景的扩展性。

实验验证表明，系统能够稳定完成事件解析、数据增强、城市与会员维度聚合、风险客户自动识别、批次级 KPI 统计以及 HTML 实时仪表盘刷新，各环节输出结果通过了累计事件数、批次总收入、平均消费、折扣订单数、高价值订单数、风险客户数等 9 项指标的一致性验证。此外，系统在空批次过滤、追加型 CSV 表头去重、路径安全校验、空闲超时自动停止和异常退出资源释放等方面均通过了流程稳定性测试。本文的工作将项目从单一的批处理分析拓展为"批处理 + 实时处理"双链路系统，使分析能力从历史数据洞察延伸到最新交易动态的实时监测。

\textbf{关键词：} 电商客户行为分析；Spark Structured Streaming；实时处理；流式分析；Kafka；风险客户识别；数据可视化
\end{abstract}

\thispagestyle{empty}

\newpage
\tableofcontents
\newpage

\section{绪论}

\subsection{研究背景}

随着互联网技术、移动支付和智能物流体系的快速发展，电子商务已经成为现代商业活动中最重要的业务形态之一。电商平台在日常运营过程中会持续产生大量数据，包括用户基本信息、浏览行为、购买记录、消费金额、会员等级、优惠券使用情况、商品评价以及客户满意度等。这些数据不仅记录了用户与平台之间的交互过程，也反映了平台的经营状况、客户偏好变化和潜在业务风险。因此，如何对电商客户行为数据进行高效处理和深入分析，已经成为电商平台提升运营效率和增强用户黏性的重要问题。

在传统数据分析场景中，批处理分析是一种常见方式。批处理通常面向已经积累完成的历史数据，通过一次性读取全量数据并进行清洗、统计、建模和可视化，从而获得用户分群、消费趋势、城市分布和满意度特征等分析结果。这种方式适合处理规模较大的静态数据，能够帮助平台从整体层面理解历史业务表现。然而，电商业务具有明显的动态性和实时性特征，交易事件、客户评价和购买行为往往是持续产生的。如果仅依赖离线批处理分析，系统只能在固定时间间隔后更新结果，难以及时反映最新业务变化，也难以及时发现异常交易、客户满意度下降或潜在客户流失等问题。

与批处理不同，实时处理和流式分析更加关注“正在发生”的业务状态。通过对持续到达的数据流进行监听和处理，系统可以在较短时间内完成数据解析、指标计算、风险识别和结果更新。例如，当平台持续产生新的交易事件时，实时处理模块可以动态统计当前批次的交易数量、总收入、平均消费金额和高价值订单数量；当某些客户出现长时间未购买、评分偏低或满意度下降等情况时，系统也可以及时将其识别为潜在风险客户。由此可见，实时处理不仅可以提升数据分析的时效性，也可以为业务运营提供更及时的决策依据。

在电商业务中，实时 KPI 监控、风险客户识别和会员行为分析都具有较高的应用价值。实时 KPI 可以帮助运营人员快速掌握当前业务运行情况，例如最新交易规模、收入变化、折扣订单数量和高价值订单数量；风险客户识别可以帮助平台提前发现可能流失的用户，为后续的精准营销、优惠推荐或客户关怀提供依据；会员行为分析则可以揭示不同会员等级客户在消费能力、评分水平和订单价值方面的差异，为会员体系优化和用户分层运营提供支持。因此，构建一套稳定、可扩展且易于展示的实时处理与流式分析模块，对于完善电商客户行为分析系统具有重要意义。

本项目面向电商客户行为分析场景，设计并实现了一个实时处理与流式分析模块。该模块以原始电商客户行为数据为基础，通过事件生成器模拟平台持续产生交易事件的过程，并使用 Spark Structured Streaming 对实时事件流进行微批次处理。系统在处理过程中完成事件留存、数据增强、城市维度聚合、会员维度聚合、风险客户识别和 KPI 汇总，并进一步生成 CSV 结果文件、Markdown 摘要以及 HTML 实时仪表盘。通过该模块，系统不仅能够分析历史数据，还能够对最新交易动态进行持续监测，从而增强整个电商客户行为分析系统的实时性和实用性。

\subsection{研究意义}

从技术角度看，实时处理模块体现了大数据系统中流式计算的核心思想。传统批处理通常采用“先存储、后计算”的方式，适合对已经落盘的历史数据进行集中处理；而流式计算则更强调对持续到达的数据进行增量处理，要求系统能够在数据产生后尽快完成解析、计算和输出。本项目基于 Spark Structured Streaming 构建实时处理链路，将实时交易事件抽象为不断增长的数据流，并通过微批次机制周期性触发处理逻辑。这种设计既保留了 Spark DataFrame API 在表达能力和开发效率方面的优势，又能够适配实时数据分析场景，有助于理解现代大数据平台中批流结合的系统设计思想。

从业务角度看，实时处理模块能够为电商平台提供更及时的运营决策支持。电商平台的用户行为变化往往具有较强的即时性，例如用户短期内评分下降、长时间未再次购买、消费金额异常变化等，都可能预示着客户满意度降低或流失风险上升。如果系统只能依赖离线批处理结果，运营人员往往只能在事后观察到这些变化，难以及时采取干预措施。通过实时处理模块，平台可以在新交易事件到达后快速更新关键指标，并识别风险客户，从而为客户维护、营销活动调整和会员运营策略优化提供依据。

从系统建设角度看，实时处理链路与批处理链路之间具有明显的互补关系。批处理分析侧重于对全量历史数据进行整体性分析，适合回答“过去发生了什么”以及“长期趋势如何”等问题；实时处理分析则侧重于对最新事件进行动态监测，适合回答“当前正在发生什么”以及“是否出现新的风险或异常”等问题。二者结合后，系统既能够从历史数据中挖掘客户行为规律，也能够对最新业务变化做出及时响应，从而形成更加完整的数据分析能力。

此外，本项目中的实时处理模块还具有一定的工程实践意义。系统同时支持文件流模式和 Kafka 模式：文件流模式不依赖额外消息队列，便于在教学和实验环境中快速运行、调试和展示；Kafka 模式则为后续扩展到更接近生产环境的分布式实时处理架构提供了基础。通过这种双模式设计，系统在保持实验可复现性的同时，也具备了一定的可扩展性和工程迁移价值。

\subsection{本文主要工作}

本文围绕电商客户行为分析系统中的实时处理与流式分析模块展开设计、实现和验证，主要工作包括以下几个方面。

第一，设计并实现实时事件生成器。为了模拟电商平台中持续产生交易数据的场景，本文没有采用完全随机生成数据的方式，而是以原始电商客户行为数据为种子，对消费金额、购买件数、平均评分、距上次购买天数、折扣使用情况和年龄等字段进行受控随机扰动。通过这种方式生成的交易事件既保留了原始数据集的统计特征，又能够体现实时数据的动态变化。同时，系统为每条事件生成唯一事件编号和事件时间，并根据评分和购买间隔动态推导客户满意度等级，从而保证事件数据在字段含义和业务逻辑上的一致性。

第二，基于 Spark Structured Streaming 构建流处理引擎。本文使用 Spark 的结构化流处理机制监听输入目录或 Kafka Topic，并通过统一的数据 Schema 将实时事件解析为结构化 DataFrame。系统采用微批次处理方式，在每个批次到达时触发完整的数据处理流程，包括空批次过滤、事件排序、数据增强、多维聚合、风险客户筛选、KPI 统计和结果输出。通过 \texttt{foreachBatch} 机制，系统能够在每个微批次内复用较为成熟的批处理 DataFrame 操作，从而降低流式计算逻辑的实现复杂度。

第三，设计风险客户识别逻辑和高价值订单识别逻辑。在风险客户识别方面，本文综合考虑客户距上次购买天数、满意度等级和平均评分等因素，将长时间未购买且满意度较低或评分较低的客户标记为潜在风险客户。该规则能够从时间维度和情感维度同时刻画客户流失风险。在高价值订单识别方面，系统将消费金额达到一定阈值的订单标记为高价值订单，用于后续 KPI 统计和会员行为分析。这两类派生指标增强了原始事件数据的业务表达能力。

第四，实现城市维度、会员维度和批次 KPI 的实时聚合统计。城市维度分析主要统计不同城市的事件数量、平均消费金额和购买商品总件数，用于观察区域消费差异；会员维度分析主要统计不同会员等级的事件数量、平均消费、平均评分和高价值订单数量，用于比较不同会员群体的行为特征；批次 KPI 统计则记录每个微批次的事件数、累计事件数、总收入、平均消费、平均评分、折扣订单数、高价值订单数和风险客户数等核心指标，用于形成实时处理过程中的时间序列记录。

第五，设计并生成多种形式的实时分析输出。系统将处理结果保存为多类 CSV 文件，包括最新批次事件、全量事件历史、城市聚合结果、会员聚合结果、风险客户明细、批次 KPI 历史和最新批次指标等。其中，快照类文件采用覆盖写入方式，以便仪表盘读取最新结果；历史类文件采用追加写入方式，以便保留处理过程中的完整记录。在此基础上，系统进一步生成 Markdown 摘要和 HTML 实时仪表盘，使分析结果能够以文本和可视化页面的形式进行展示。

第六，预留 Kafka 模式以增强系统扩展性。除默认文件流模式外，本文还设计了基于 Kafka 的事件生产和消费方式。事件生成器可以将 JSON 格式交易事件发送至 Kafka Topic，流处理引擎也可以从 Kafka 中读取消息并解析为结构化数据。虽然在教学演示场景中默认采用文件流模式以降低环境复杂度，但 Kafka 模式的保留使系统具备进一步扩展到分布式实时处理环境的能力，也使该模块更接近实际生产系统中的消息队列式架构。

综上，本文的主要工作不仅包括实时数据处理链路的功能实现，也包括对系统架构、数据生成方式、流式计算逻辑、业务指标设计、结果输出形式和扩展能力的综合设计。通过这些工作，本文构建了一个能够支撑电商客户行为动态监测的实时处理模块，为后续进一步扩展窗口分析、实时告警和客户流失预测等功能奠定了基础。


\section{系统总体设计}

\subsection{系统目标}

本系统的实时处理与流式分析模块面向电商客户行为分析场景，旨在构建一条从交易事件生成、实时数据接入、流式计算处理到结果展示输出的完整数据链路。该模块并不局限于对静态历史数据进行分析，而是进一步模拟电商平台在运行过程中持续产生交易事件的业务场景，使系统能够对最新客户行为和交易状态进行动态监测。通过实时处理模块，电商客户行为分析系统可以从单纯的离线统计扩展为同时具备历史分析和实时监控能力的数据分析平台。

在业务场景模拟方面，系统需要尽可能接近真实电商平台中交易事件持续产生的过程。真实电商业务中，用户购买商品、使用折扣、提交评分、产生消费金额等行为并不是一次性完成的，而是随着时间不断发生。为了在实验环境中模拟这一过程，本系统设计了实时事件生成器。该生成器以原始电商客户行为数据为基础，通过对已有记录进行受控随机扰动，生成新的交易事件。这样既避免了完全随机生成数据导致业务含义不合理的问题，也能够保留原始数据集中的客户行为特征，使生成的数据更加接近真实业务场景。

在实时数据处理方面，系统需要支持对最新交易事件的持续监听和微批次处理。传统批处理方式通常在数据完全准备好之后再进行集中计算，适合分析历史数据，但难以及时反映最新业务变化。实时处理模块则要求系统能够不断接收新增事件，并在较短时间内完成数据解析、统计分析和结果更新。因此，本系统采用 Spark Structured Streaming 作为流处理框架，将持续写入的 JSON 文件或 Kafka 消息抽象为实时数据流，并通过微批次方式对新增数据进行周期性计算。每个微批次都会触发一次完整的数据处理流程，从而保证分析结果能够随着新事件的到达而不断更新。

在结果输出方面，系统需要保证分析结果具备可解释性、可追踪性和可展示性。可解释性体现在系统输出的各类指标均具有明确业务含义，例如批次收入、平均消费、风险客户数量、高价值订单数量等；可追踪性体现在系统会保存批次级事件文件、累计事件历史和 KPI 历史记录，便于对运行过程进行复核；可展示性则体现在系统会根据处理结果生成 HTML 实时仪表盘和 Markdown 摘要，使用户能够直观查看实时处理状态。通过结构化 CSV 文件与可视化页面相结合的方式，系统既方便后续数据分析，也便于课程展示和结果说明。

在运行环境适配方面，系统需要同时兼顾教学演示的简洁性和生产级架构的可扩展性。课程实验或单机演示通常希望降低外部依赖，避免复杂的 Kafka 和 Zookeeper 配置，因此系统默认支持文件流模式，即事件生成器将 JSON Lines 文件写入监听目录，由 Spark Structured Streaming 自动读取新增文件并触发处理。同时，系统也预留 Kafka 模式，使事件生成器可以将交易事件发送至 Kafka Topic，流处理引擎从 Kafka 中消费消息并完成实时分析。通过这种双模式设计，系统既能够在本地环境中快速运行和验证，也具备进一步扩展到分布式实时处理架构的基础。

总体而言，本系统的设计目标可以概括为：以电商客户行为数据为基础，模拟真实交易事件的持续产生过程，利用 Spark Structured Streaming 实现实时数据接入和微批次处理，并通过多种输出形式展示实时分析结果。该模块不仅补充了批处理分析在时效性方面的不足，也为后续扩展实时告警、窗口统计、客户流失预测和生产级消息队列架构提供了基础。

\subsection{实时处理链路总体架构}

本系统的实时处理链路采用模块化设计思想，将完整的实时分析流程拆分为事件生成、流式处理和结果展示三个相对独立的功能层。三个功能层分别由 \texttt{realtime\_producer.py}、\texttt{realtime\_processor.py} 和 \texttt{realtime\_reporting.py} 三个核心脚本实现。各脚本之间通过文件目录、Kafka Topic 和输出结果文件进行衔接，形成一条较为清晰的数据流水线。

\begin{table}[H]
\centering
\caption{实时处理模块核心脚本}
\begin{tabular}{llll}
\toprule
脚本名称 & 模块角色 & 输入 & 输出 \\
\midrule
\texttt{realtime\_producer.py} & 事件生成器 & 原始 CSV 种子数据 & JSON 事件文件 / Kafka 消息 \\
\texttt{realtime\_processor.py} & 流处理引擎 & JSON 文件流 / Kafka Topic & CSV 结果与仪表盘 \\
\texttt{realtime\_reporting.py} & 报表生成器 & 处理结果 CSV & HTML 仪表盘与 Markdown 摘要 \\
\bottomrule
\end{tabular}
\end{table}

事件生成层由 \texttt{realtime\_producer.py} 实现，位于实时处理链路的起点。该脚本负责从原始电商客户行为数据中读取种子记录，并基于这些记录生成模拟交易事件。每条事件包含客户编号、性别、年龄、城市、会员类型、消费金额、购买数量、平均评分、折扣使用情况、距上次购买天数、满意度等级、事件编号和事件时间等字段。为了保证生成事件的真实性，事件生成器不会脱离原始数据随意构造字段，而是在原始记录基础上施加合理扰动，使事件既具有随机变化，又保持业务逻辑上的合理性。在文件流模式下，事件生成器会将事件按批次写入 JSON Lines 文件；在 Kafka 模式下，则会将事件作为 JSON 消息发送到指定 Kafka Topic。

流式处理层由 \texttt{realtime\_processor.py} 实现，是整个实时处理模块的核心部分。该脚本基于 Spark Structured Streaming 构建，可以根据运行参数选择从文件目录读取 JSON 事件流，或者从 Kafka Topic 中消费实时消息。无论数据来自文件流还是 Kafka，系统都会将其解析为统一 Schema 的结构化 DataFrame，并在此基础上执行后续处理逻辑。处理内容包括事件排序、事件保存、数据增强、城市维度聚合、会员维度聚合、风险客户识别和批次 KPI 统计等。通过 \texttt{foreachBatch} 机制，流处理引擎能够在每个微批次到达时执行一套完整的分析流程，从而实现对实时事件的连续处理。

结果展示层由 \texttt{realtime\_reporting.py} 实现，位于实时处理链路的末端。该脚本主要负责读取流处理阶段生成的 CSV 结果文件，并将其渲染为更适合展示和阅读的形式。系统输出的 HTML 仪表盘包含核心 KPI 指标卡片、城市平均消费展示、会员类型事件分布、批次 KPI 历史、风险客户列表和最新事件明细等内容。Markdown 摘要则以文本形式概括当前实时处理任务的运行情况，包括处理批次数、累计事件数、最新批次收入、平均消费和风险客户数量等。通过报表生成器，系统能够将底层结构化数据转化为直观的分析结果，提高报告展示效果和结果可读性。

从整体结构看，三个核心脚本之间形成了较清晰的职责边界。事件生成器只负责产生实时事件，不参与后续分析；流处理引擎只负责对输入事件进行计算和输出结果，不关心页面样式和展示细节；报表生成器只负责读取结果并生成仪表盘，不影响事件生成和流式计算过程。这种松耦合结构提高了系统的可维护性和可扩展性。当后续需要更换数据源、接入数据库或改进仪表盘样式时，可以针对某一模块进行调整，而不必重写整个实时处理链路。

系统整体架构可以概括为如下数据流：

\begin{center}
\texttt{原始 CSV 数据}
$\rightarrow$
\texttt{实时事件生成器}
$\rightarrow$
\texttt{JSON 文件流 / Kafka Topic}
$\rightarrow$
\texttt{Spark 流处理引擎}
$\rightarrow$
\texttt{CSV 分析结果}
$\rightarrow$
\texttt{HTML 仪表盘与 Markdown 摘要}
\end{center}

该架构实现了从原始历史数据到实时事件、从实时事件到流式分析结果、从分析结果到可视化报告的完整转换过程。系统既能以较低成本完成本地演示，也能够在 Kafka 模式下扩展为更接近生产环境的实时处理链路。

\begin{figure}[H]
\centering
\includegraphics[width=\textwidth]{figures/pic1.png}
\caption{实时处理链路总体架构图：从原始 CSV 数据经事件生成器、文件流/Kafka 双模式输入、Spark Structured Streaming 流处理引擎到 CSV 结果与 HTML 仪表盘输出的完整数据流转过程}
\label{fig:architecture}
\end{figure}

\subsection{数据流转过程}

本系统的数据流转过程贯穿事件生成、事件接入、实时计算和结果展示四个主要环节。各环节之间通过标准化的数据格式进行衔接，使系统能够保持较好的可读性和可维护性。

在种子数据读取环节，事件生成器从原始电商客户行为 CSV 文件中读取客户记录。原始数据中包含客户的基础属性和行为属性，例如客户编号、性别、年龄、城市、会员类型、消费金额、购买商品数量、平均评分、折扣使用情况、距上次购买天数和满意度等级等。这些字段为实时事件生成提供了基础。由于种子数据来源于原始数据集，因此生成的实时事件能够延续原始数据的分布特征，避免出现明显偏离业务场景的异常样本。

在实时事件生成环节，系统从种子数据池中抽取记录，并对关键字段进行受控随机扰动。消费金额会在原值附近浮动，购买件数会在合理范围内变化，平均评分会被限制在指定区间内，距上次购买天数也会根据扰动规则进行更新。与此同时，系统会为每条事件分配唯一事件编号，并记录事件生成时间。满意度等级则根据变异后的评分和购买间隔重新推导，而不是直接沿用原始记录中的字段值。通过这种方式，每条生成事件都能够保持内部逻辑一致，即评分、购买间隔和满意度之间存在合理关系。

在事件写入环节，系统根据运行模式选择不同的数据输出方式。文件流模式下，事件生成器按照批次将交易事件写入监听目录，每个批次对应一个 JSON Lines 文件。JSON Lines 文件中每一行表示一条完整事件，适合用于模拟实时文件流输入。Kafka 模式下，事件生成器将每条交易事件序列化为 JSON 消息，并发送到指定 Kafka Topic。虽然两种模式的数据传输方式不同，但事件本身采用统一结构，因此后续处理过程不需要针对业务字段进行重复适配。

在流式数据监听环节，Spark Structured Streaming 持续监听输入源。对于文件流模式，Spark 会监控输入目录中新出现的 JSON 文件，并将新增文件中的记录读入流式 DataFrame；对于 Kafka 模式，Spark 会持续消费 Kafka Topic 中的新消息，并将消息内容解析为结构化数据。系统通过统一 Schema 对输入事件进行约束，确保不同数据源进入处理阶段后具有一致的数据结构。这种设计使文件流模式和 Kafka 模式能够共用同一套数据增强和聚合逻辑。

在微批次分析环节，流处理引擎按照设定的触发间隔对新增事件进行批次化处理。每个微批次到达后，系统会检查该批次是否包含有效事件。对于非空批次，系统会将事件按照事件时间和事件编号排序，并分别写入批次归档文件、最新事件快照文件和全量事件历史文件。随后，系统在原始字段基础上新增风险客户标记和高价值订单标记，并基于增强后的数据计算城市维度统计、会员维度统计、风险客户明细和批次 KPI 指标。通过这种处理方式，系统能够在每个批次中同时完成明细保存、统计聚合和业务识别。

在结果展示环节，系统根据流处理输出的 CSV 文件刷新 HTML 仪表盘和 Markdown 摘要。HTML 仪表盘主要用于展示实时处理结果，包括累计事件数、最新批次收入、风险客户数量、城市消费表现、会员等级分布、风险客户列表和最新交易事件等。Markdown 摘要则提供较为简洁的文本化运行概览，便于在报告或实验记录中引用。由于仪表盘和摘要会在每个批次处理完成后更新，因此用户可以通过查看输出文件实时了解系统运行状态。

整体来看，系统数据流转过程具有较强的连续性和可追踪性。事件从原始 CSV 数据中产生，经由文件目录或 Kafka Topic 进入流处理引擎，再被加工为多类结构化结果和可视化报告。每个环节都有明确输入和输出，既方便调试，也便于验证结果正确性。这种设计使实时处理模块能够较好地适配课程项目中的演示需求，同时也为未来接入真实数据源和生产级消息队列提供了扩展空间。

\subsection{文件流模式与 Kafka 模式}

本系统在实时数据接入层设计了文件流模式和 Kafka 模式两种运行方式。两种模式面向不同运行场景，但在事件结构、处理逻辑和输出结果方面保持一致。这样的设计既降低了实验环境中的运行门槛，也保留了向生产级实时处理架构扩展的可能性。

\begin{table}[H]
\centering
\caption{两种实时处理模式对比}
\begin{tabular}{lll}
\toprule
对比维度 & 文件流模式 & Kafka 模式 \\
\midrule
运行参数 & \texttt{--mode files} & \texttt{--mode kafka} \\
输入输出方式 & JSON 文件目录 & Kafka Topic \\
环境依赖 & 低，无需额外服务 & 需要 Kafka 与 Zookeeper \\
适用场景 & 教学演示、单机实验 & 分布式、生产级实时处理 \\
优点 & 简单、可复现、易调试 & 扩展性强、接近真实生产环境 \\
\bottomrule
\end{tabular}
\end{table}

文件流模式是系统默认采用的实时处理方式。在该模式下，事件生成器将每个批次的交易事件写入本地监听目录，流处理引擎通过 Spark Structured Streaming 读取该目录中的新增 JSON 文件。由于文件流模式不需要额外部署 Kafka、Zookeeper 等服务，运行环境较为简单，适合课程项目展示、单机实验和功能验证。使用者可以直接观察输入目录中新生成的批次文件，以及输出目录中实时更新的 CSV 文件和 HTML 仪表盘，从而直观理解实时处理链路的运行过程。

文件流模式的另一个优势是可复现性较强。事件生成器可以通过固定随机种子和批次参数控制生成结果，使实验运行过程更容易复核。在调试阶段，文件流模式也更便于定位问题。例如，当某个批次处理结果异常时，可以直接查看对应的 JSON 输入文件和批次输出 CSV 文件，判断问题出现在事件生成、数据解析还是聚合统计环节。因此，在教学演示和本地开发阶段，文件流模式能够显著降低系统运行和调试成本。

Kafka 模式面向更接近生产环境的实时数据处理场景。在实际电商系统中，交易事件通常由前端服务、订单服务或日志采集系统持续产生，并通过消息队列传递给下游计算模块。Kafka 作为常见的分布式消息系统，能够支持高吞吐量、可扩展和解耦的数据传输方式。本项目预留 Kafka 模式，使事件生成器可以将 JSON 格式的交易事件发送到 Kafka Topic，流处理引擎则从该 Topic 中消费消息，并将消息解析为结构化事件数据。

Kafka 模式的价值主要体现在系统扩展能力上。当事件规模增大或需要多个服务同时消费同一数据流时，Kafka 能够提供比本地文件目录更强的数据分发能力和容错能力。通过 Kafka Topic，事件生产和事件消费之间可以解耦，生产端只负责写入消息，消费端可以根据自身处理能力持续读取消息。对于未来可能扩展的实时告警、在线推荐、客户流失预测等模块，Kafka 模式能够提供更适合的事件接入基础。

虽然文件流模式和 Kafka 模式在输入机制上存在差异，但系统在处理层面保持了统一设计。文件流模式直接读取 JSON 文件并应用统一 Schema；Kafka 模式则从消息中提取 JSON 字符串，再使用相同 Schema 解析为结构化 DataFrame。完成解析后，两种模式都会进入相同的数据增强、聚合统计、风险识别和结果输出流程。这样的设计减少了重复代码，也保证了不同运行模式下分析结果的一致性。

从本项目的实际需求来看，文件流模式更适合当前课程实验环境，能够以较低配置成本完成完整实时处理链路的演示；Kafka 模式则体现了系统面向真实业务环境的扩展能力。双模式设计使系统既能满足“能够运行、便于展示、结果可验证”的实验要求，也具备“可扩展、可迁移、接近生产架构”的工程特征。通过这种设计，本实时处理模块在教学可用性和工程合理性之间取得了较好的平衡。

\section{事件生成器设计与实现}

\subsection{设计目标}

事件生成器是实时处理链路的起点，其主要作用是模拟电商平台在运行过程中持续产生交易事件的场景。在真实电商业务中，用户的购买、评分、折扣使用和会员消费等行为会随着时间不断发生，这些行为构成了实时数据流的来源。由于本项目使用的是静态电商客户行为数据集，若要验证实时处理链路的有效性，就需要在实验环境中构造出持续到达的新交易事件。因此，\texttt{realtime\_producer.py} 被设计为一个事件生成模块，用于将原始静态数据转化为可被流处理引擎持续消费的模拟实时事件。

在数据真实性方面，事件生成器需要避免完全随机生成数据。若直接随机产生客户编号、消费金额、评分、会员类型和满意度等级等字段，虽然可以构造出形式上完整的事件，但这些事件很容易偏离真实业务场景。例如，随机生成的高消费客户可能与会员等级不匹配，评分和满意度之间可能缺乏逻辑关系，距上次购买天数也可能与客户满意度不一致。此类数据虽然能够作为程序输入，却难以支撑有意义的业务分析。因此，本系统采用基于原始数据的种子变异方法，将真实数据记录作为事件生成基础，在保持原有统计特征的同时引入适度随机变化。

在动态性方面，事件生成器需要让数据能够体现实时业务中的变化过程。原始数据集中的记录是固定的，直接重复读取这些记录无法体现新事件产生的随机性和时间变化。为此，系统在每次生成事件时都会对消费金额、购买件数、平均评分、距上次购买天数、折扣使用情况和年龄等字段进行受控扰动。通过这种方式，生成事件既不是对原始记录的简单复制，也不是脱离业务背景的随机样本，而是在原始客户行为基础上形成的合理变化版本。

在字段完整性和逻辑一致性方面，事件生成器需要保证每条输出事件都能够被后续流处理模块正确解析和使用。每条事件不仅要包含客户编号、城市、会员类型、消费金额、购买数量、评分等业务字段，还需要包含事件唯一标识和事件时间等流式处理所需字段。同时，系统会根据变异后的评分和购买间隔重新推导满意度等级，使满意度字段与客户行为特征之间保持一致。通过这种设计，事件生成器输出的数据能够直接进入 Spark Structured Streaming 处理链路，而不需要额外的数据修复或字段补全。

在实验可复现性方面，事件生成器需要保证相同参数下能够得到稳定结果。课程项目中的实验结果通常需要反复运行、调试和验证，如果每次运行产生的数据差异过大，就会增加结果比对和问题排查的难度。因此，系统在事件生成过程中使用固定随机种子，并允许通过批次大小、生成批次数和批次间隔等参数控制输出过程。这样既能够模拟流式数据持续到达的效果，也能够在需要时复现实验结果，便于后续对实时处理链路进行验证。

综上，事件生成器的设计目标可以概括为：以原始电商客户行为数据为基础，构造持续到达、字段完整、逻辑自洽、具有统计真实性并且结果可复现的模拟交易事件，为后续流处理引擎提供稳定可靠的实时输入数据。

\subsection{种子数据加载}

事件生成器在运行开始时，会从原始 CSV 数据集中读取全部客户行为记录，并将这些记录存入内存中的种子池。种子池是后续事件生成的基础，每次生成新事件时，系统都会从种子池中抽取一条原始记录，并在该记录的基础上进行字段扰动。与完全随机生成方式相比，基于种子数据的生成方法能够更好地保留原始数据中的客户结构、城市分布、会员等级分布和消费行为特征，使模拟事件更加符合电商业务场景。

在 CSV 文件读取方面，系统使用 Python 内置的 \texttt{csv.DictReader} 读取原始数据。该读取方式会将 CSV 文件中的每一行解析为字典结构，表头字段作为字典键，对应单元格内容作为字典值。相比按位置读取列表形式数据，字典形式具有更好的可读性和可维护性。例如，后续处理可以直接通过 \texttt{row["total\_spend"]}、\texttt{row["average\_rating"]} 或 \texttt{row["membership\_type"]} 访问字段，而不需要依赖字段所在列的位置。这种方式降低了字段顺序变化对程序的影响，也使事件生成逻辑更加清晰。

在编码处理方面，系统以 \texttt{utf-8-sig} 编码打开 CSV 文件。该编码能够兼容带有 BOM 标记的 UTF-8 文件，避免表头字段在读取时出现隐藏字符。例如，如果 CSV 文件开头包含 BOM，而程序直接使用普通 \texttt{utf-8} 读取，则第一个字段名可能会带有不可见字符，导致后续通过字段名访问数据时出现错误。使用 \texttt{utf-8-sig} 可以自动处理这一问题，提高数据读取过程的稳健性。

在空数据异常检查方面，系统会在读取完成后判断记录列表是否为空。如果原始 CSV 文件不存在有效记录，事件生成器会抛出异常并终止运行，而不是继续执行后续生成逻辑。该检查能够避免在种子池为空时进行随机抽样，防止程序在运行过程中出现难以理解的错误。对于实时处理链路而言，输入数据是后续所有处理步骤的基础，因此在事件生成阶段提前进行数据有效性检查具有必要性。

种子池在事件生成过程中发挥着核心作用。每一条实时事件都来源于种子池中的一条原始客户记录，并在该记录基础上进行合理扰动。这样生成的数据能够保留原始客户行为数据中的总体分布特征。例如，如果原始数据中某些城市客户数量较多，生成事件中这些城市也更可能频繁出现；如果某些会员类型在原始数据中具有较高消费金额，变异后的事件也会在一定程度上延续这一特征。由此可见，种子池不仅提供了事件生成所需的字段模板，也保证了模拟数据在统计意义上的合理性。

种子数据加载过程可以概括为如下逻辑：

\begin{verbatim}
def load_seed_rows(path):
    with open(path, "r", encoding="utf-8-sig", newline="") as handle:
        reader = csv.DictReader(handle)
        rows = list(reader)
    if not rows:
        raise ValueError(f"No rows found in source file: {path}")
    return rows
\end{verbatim}

该函数完成了文件打开、CSV 解析、记录列表构造和空数据检查四个步骤。返回的 \texttt{rows} 即为后续事件生成所使用的种子池。由于种子池在程序启动时一次性加载，后续生成事件时可以直接在内存中抽取记录，因此能够满足本项目模拟实时数据生成的性能需求。

\subsection{受控数据变异机制}

事件生成器的核心逻辑是对原始客户记录进行受控数据变异。所谓受控变异，是指系统并不随意修改字段，而是在预设范围内对部分数值型字段和布尔型字段进行随机扰动，同时保留城市、会员类型等基础属性字段的业务语义。通过这种方式，生成事件能够体现实时数据的变化，又不会破坏原始数据的合理分布。

在本项目中，受控数据变异主要由 \texttt{mutate\_row()} 函数实现。该函数接收一条原始客户记录和事件编号作为输入，输出一条新的交易事件。变异过程涉及消费金额、购买件数、平均评分、距上次购买天数、折扣使用情况和年龄等字段。这些字段之所以被选为扰动对象，是因为它们在真实业务中本身具有动态变化特征。例如，同一客户在不同时间的消费金额可能不同，购买商品数量可能不同，对平台或商品的评分也可能发生变化。

\begin{table}[H]
\centering
\caption{交易事件字段变异规则}
\begin{tabular}{lll}
\toprule
字段 & 变异方式 & 约束条件 \\
\midrule
\texttt{total\_spend} & 原值加随机扰动 & 下限为 120.0 \\
\texttt{items\_purchased} & 原值加整数扰动 & 下限为 1 \\
\texttt{average\_rating} & 原值加评分扰动 & 限制在 2.8 至 5.0 \\
\texttt{days\_since\_last\_purchase} & 原值加天数扰动 & 下限为 1 \\
\texttt{discount\_applied} & 按概率翻转布尔值 & 保持布尔类型 \\
\texttt{age} & 原值加小幅扰动 & 限制在 18 至 70 \\
\bottomrule
\end{tabular}
\end{table}

在消费金额变异方面，系统以原始 \texttt{total\_spend} 为基础，增加一个随机浮动值，并设置最低消费金额下限。这样可以模拟同一类客户在不同交易中的消费波动，同时避免出现过低或不合理的消费金额。消费金额是后续高价值订单识别、城市平均消费统计和批次收入计算的重要字段，因此该字段的合理性会直接影响实时分析结果的可信度。

在购买件数变异方面，系统对 \texttt{items\_purchased} 进行整数扰动，并保证结果不低于 1。购买件数是订单行为中的基本指标，不应出现 0 或负数。通过设置下限，系统保证每条生成事件都代表一次有效购买行为。该字段后续会参与城市维度的购买总件数统计，因此保持其数值合理性具有实际意义。

在平均评分变异方面，系统对 \texttt{average\_rating} 进行小幅浮动，并将结果限制在合理区间内。评分通常具有明确上下界，如果扰动后评分超过正常范围，会导致满意度推导和风险客户识别失真。因此，系统将评分限制在 2.8 至 5.0 之间，使评分既可以体现客户体验变化，又不会脱离常见评分范围。该字段后续不仅参与平均评分统计，也会影响风险客户识别。

在距上次购买天数变异方面，系统基于原始 \texttt{days\_since\_last\_purchase} 增加随机天数扰动，并保证结果不低于 1。该字段反映客户最近一次购买与当前事件之间的时间间隔，是判断客户活跃程度和流失风险的重要指标。通过扰动该字段，系统可以模拟客户购买频率变化，使风险客户识别逻辑能够在实时事件中发挥作用。

在折扣使用情况变异方面，系统不会每次都改变 \texttt{discount\_applied} 字段，而是按照一定概率翻转其布尔值。该设计符合真实业务中用户是否使用折扣存在不确定性的特点。如果折扣字段完全沿用原始记录，则实时事件变化不足；如果每次都随机重置，则可能破坏原始客户行为特征。按概率翻转能够在稳定性和变化性之间取得平衡。

在年龄变异方面，系统只允许小幅变化，并将年龄限制在 18 至 70 之间。虽然客户年龄在真实业务中不会在短时间内频繁变化，但在模拟数据生成场景中，小幅扰动可以用于扩展样本变化，同时不会改变客户群体的基本结构。年龄字段不作为本模块实时 KPI 的核心指标，但仍然属于客户基础属性，因此需要保持合理范围。

除上述字段外，事件生成器还会为每条事件增加 \texttt{event\_id} 和 \texttt{event\_time} 字段。\texttt{event\_id} 用于唯一标识事件，通常采用按序递增的编号形式，例如 \texttt{evt-000001}；\texttt{event\_time} 用于记录事件生成时间，通常采用 ISO 8601 格式表示。二者共同增强了事件的可追踪性，使后续流处理模块能够按照事件时间和事件编号对数据进行排序和归档。

受控数据变异机制的意义在于，它在真实数据和模拟数据之间建立了过渡关系。生成事件既继承原始数据的基本特征，又通过扰动体现实时业务变化；既能满足程序测试和系统演示需求，又能为后续风险识别、KPI 统计和可视化展示提供具有业务含义的数据基础。

\subsection{满意度等级推导}

满意度等级是电商客户行为分析中的重要业务字段，能够反映客户对平台、商品或服务的总体体验。在事件生成器中，系统并不是直接随机赋予满意度等级，也不是简单沿用原始记录中的满意度字段，而是根据变异后的平均评分和距上次购买天数进行动态推导。这样可以保证满意度等级与当前事件中的客户行为特征保持一致。

系统采用的满意度推导逻辑可以概括为：当客户平均评分较高且距离上次购买时间较短时，说明客户近期仍保持较高活跃度，并且对平台体验较好，因此将其判定为 \texttt{Satisfied}；当客户评分较低，或者已经较长时间没有购买行为时，说明客户可能存在体验不佳或活跃度下降的问题，因此将其判定为 \texttt{Unsatisfied}；介于两者之间的情况则被判定为 \texttt{Neutral}。

具体逻辑可以表示为：

\begin{verbatim}
def derive_satisfaction(average_rating, days_since_last_purchase):
    if average_rating >= 4.3 and days_since_last_purchase <= 28:
        return "Satisfied"
    if average_rating < 3.7 or days_since_last_purchase > 38:
        return "Unsatisfied"
    return "Neutral"
\end{verbatim}

该规则综合考虑了情感维度和时间维度。平均评分代表客户对平台或服务的主观评价，距上次购买天数则代表客户近期活跃程度。单独使用评分可能无法识别长时间未购买但评分尚可的客户，单独使用购买间隔也可能误判短期未购买但满意度较高的客户。将两类因素结合起来，可以更合理地刻画客户当前状态。

从业务逻辑角度看，高评分和近期购买通常意味着客户对平台仍保持较强认可，因此被归入满意客户；低评分往往直接反映客户体验问题，因此即使购买间隔不长，也可能被归入不满意客户；长时间未购买则可能表明客户活跃度下降，即使其历史评分不低，也需要被关注。中性客户则代表状态相对稳定但不够突出的一类客户，既没有表现出明显满意，也没有达到高风险或强不满状态。

满意度推导规则还为后续风险客户识别提供了基础。在流处理阶段，系统会结合满意度等级、平均评分和距上次购买天数判断客户是否存在流失风险。如果满意度字段与评分和购买间隔之间缺乏逻辑关系，风险客户识别结果就可能出现矛盾。例如，一个评分很低且长时间未购买的客户却被标记为满意，会削弱风险识别的可信度。通过在事件生成阶段动态推导满意度，可以减少这类不一致问题，提高后续分析结果的合理性。

此外，满意度推导逻辑也增强了模拟事件的可解释性。每条事件的满意度等级都可以通过评分和购买间隔解释，而不是不可追溯的随机标签。这对于课程报告和实验展示具有重要作用，因为分析结果不仅需要程序上正确，也需要能够从业务角度说明其形成原因。

\subsection{批量发布机制}

为了模拟实时数据持续到达的过程，事件生成器采用批量发布机制。批量发布并不是一次性生成全部事件，而是按照设定的批次大小、批次数量和批次间隔逐步输出数据。该设计能够使下游 Spark Structured Streaming 处理器在运行过程中不断接收到新数据，从而触发多个微批次处理过程，较好地模拟真实流式数据分析场景。

在批次大小方面，系统允许通过参数控制每个批次生成的事件数量。默认情况下，每个批次会生成固定数量的交易事件。批次大小决定了每次流处理引擎能够读取到的数据规模。如果批次过小，处理结果可能波动较大，聚合统计意义不明显；如果批次过大，则会降低实时模拟的细粒度。因此，在实验场景中选择适中的批次大小，可以在计算效率和实时效果之间取得平衡。

在批次数量方面，系统通过参数控制事件生成器总共发布多少个批次。该参数适合课程演示和实验验证，因为实验通常不需要无限运行的数据流，而是需要在有限时间内观察完整的数据生成、流式处理和结果输出过程。通过设定固定批次数，系统能够在生成指定数量事件后自动结束，便于复现实验流程和保存结果。

在批次间隔方面，事件生成器会在相邻批次之间暂停一段时间。该时间间隔用于模拟真实业务中事件持续到达但并非一次性集中产生的特点。下游流处理引擎可以根据触发间隔发现新输入文件或新 Kafka 消息，并对其进行微批次处理。批次间隔的存在使系统能够更清晰地展示“事件生成—流式监听—批次处理—结果刷新”的动态过程。

文件流模式下，事件生成器会将每个批次写入一个独立的 JSON Lines 文件，文件名通常带有批次编号，例如 \texttt{batch\_001.json}、\texttt{batch\_002.json} 等。每个文件中包含多条 JSON 格式交易事件，每一行对应一条事件。这样的文件组织方式具有较强的可读性和可追踪性。当后续处理结果出现异常时，可以回到对应批次文件中检查原始输入事件。

Kafka 模式下，事件生成器会将每条交易事件序列化为 JSON 消息，并发送到指定 Kafka Topic。相比文件流模式，Kafka 模式不再以批次文件作为传输媒介，而是通过消息队列实现事件传递。虽然底层传输方式不同，但事件生成逻辑保持一致，生成的事件字段结构也保持一致。因此，Kafka 模式可以在不改变核心业务逻辑的基础上，将系统扩展到更接近生产环境的实时处理架构。

为了保证实验结果可复现，事件生成器还支持固定随机种子。随机种子会影响种子记录抽取、字段扰动和折扣状态翻转等随机过程。如果随机种子固定，并且输入数据和运行参数相同，则多次运行可以得到一致或高度可控的事件序列。这对于课程项目中的结果验证、截图展示和报告撰写非常重要。

批量发布机制的整体流程可以概括为：

\begin{verbatim}
for batch_id in range(1, iterations + 1):
    events = []
    for i in range(batch_size):
        seed_row = random.choice(seed_rows)
        event = mutate_row(seed_row, event_index)
        events.append(event)

    if mode == "files":
        write_json_lines(events, output_file)
    elif mode == "kafka":
        send_events_to_kafka(events, topic)

    sleep(sleep_seconds)
\end{verbatim}

该流程体现了事件生成器的基本工作方式：从种子池抽取记录，生成当前批次事件，根据运行模式写入文件或发送 Kafka 消息，然后等待下一批次生成。通过这一机制，系统能够将静态原始数据转化为连续到达的实时事件流，为后续 Spark Structured Streaming 处理提供输入基础。

综上，批量发布机制不仅实现了模拟实时数据流的功能，也增强了系统运行过程的可控性和可观察性。用户可以通过调整批次大小、批次数量和批次间隔，灵活控制实验规模和实时效果；通过文件流或 Kafka 模式选择不同的数据接入方式；通过固定随机种子保证实验结果可复现。这些设计共同保证了事件生成器在实时处理链路中的稳定性和实用性。

\begin{figure}[H]
\centering
\includegraphics[width=0.60\textwidth]{figures/pic3.png}
\caption{实时事件生成器完整工作流程：从种子数据加载、受控字段变异、满意度推导到批量发布的端到端过程}
\label{fig:event_gen_flow}
\end{figure}

\section{流处理引擎设计与实现}

\subsection{Spark Structured Streaming 简介}

流处理引擎是实时处理链路中的核心模块，负责对事件生成器持续产生的交易事件进行读取、解析、增强、聚合和输出。本系统采用 Spark Structured Streaming 作为实时计算框架，主要原因在于其能够将流式数据处理与 Spark SQL、DataFrame API 统一起来，使实时计算逻辑能够以接近批处理的方式进行表达，从而降低系统实现复杂度。

Spark Structured Streaming 的基本思想是将持续到达的数据流抽象为一张不断增长的无界表。对于文件目录或 Kafka Topic 中不断新增的数据，Spark 并不会将其视为一次性静态输入，而是将其看作持续追加到表中的新记录。每当有新数据到达时，系统会根据设置的触发间隔生成一个微批次，并在该微批次上执行预先定义好的计算逻辑。通过这种抽象方式，开发者可以使用类似处理静态 DataFrame 的方式处理流式数据，而不需要直接管理底层的事件监听、批次调度和状态更新细节。

在微批次处理机制下，流式计算任务会周期性检查输入源中是否存在新增数据。如果有新事件到达，Spark 会将这些新增数据组织成一个批次，并触发相应的处理逻辑；如果没有新事件，系统可以跳过当前批次，等待下一轮触发。微批次模式在实时性和计算稳定性之间取得了较好的平衡。相比逐条事件处理，微批次方式可以降低频繁调度带来的开销；相比传统批处理，微批次又能够在较短时间间隔内持续更新分析结果，满足近实时分析需求。

Structured Streaming 与传统 Spark DataFrame API 之间具有较强的一致性。对于静态批处理任务，开发者通常使用 \texttt{groupBy}、\texttt{agg}、\texttt{filter}、\texttt{select}、\texttt{orderBy} 等 DataFrame 操作完成数据分析。在 Structured Streaming 中，这些操作仍然可以被用于流式数据，只是输入数据由有限表变成了持续增长的无界表。本系统正是利用这一特点，在每个微批次中复用 DataFrame 操作完成城市聚合、会员聚合、风险客户筛选和 KPI 计算，使实时处理逻辑保持较高的可读性。

本项目选择 Spark Structured Streaming 还与系统整体架构有关。一方面，电商客户行为数据本身适合被表示为结构化表格数据，每条交易事件都具有固定字段，如客户编号、城市、会员类型、消费金额、评分和满意度等级等；另一方面，本系统既需要进行明细事件保存，也需要进行多维度聚合统计，这些操作都可以很好地用 DataFrame API 表达。此外，Spark Structured Streaming 同时支持文件源和 Kafka 源，能够与本项目的双模式输入设计相匹配。因此，采用 Structured Streaming 可以在保证系统功能完整的同时，减少不同数据源之间的处理差异。

在本系统中，流处理引擎主要由 \texttt{realtime\_processor.py} 实现。该脚本在启动后根据运行参数选择输入源，并对输入事件应用统一 Schema。随后，系统通过数据增强逻辑新增风险客户标记和高价值订单标记，再利用 \texttt{foreachBatch} 在每个微批次中执行完整处理流程。处理结果最终被写入 CSV 文件，并触发报表生成模块刷新 HTML 仪表盘和 Markdown 摘要。由此，Spark Structured Streaming 在本项目中承担了连接实时事件输入与实时分析结果输出的关键作用。

\subsection{Checkpoint 机制与容错设计}

在 Structured Streaming 中，Checkpoint（检查点）是保证流处理任务容错性和状态一致性的关键机制。系统会将流查询的元数据（包括输入源偏移量、已处理批次信息、状态聚合中间结果等）持久化到可靠的分布式文件系统中。当流处理任务因节点故障、主动重启或资源调度等原因中断后，重新启动的查询可以从最近的 Checkpoint 恢复处理，避免事件丢失或重复计算。

本系统为 Checkpoint 设计了专门的存储目录，默认路径为 \texttt{/opt/spark-results/realtime/checkpoints/processor/}。在流查询启动时，系统通过以下配置将 Checkpoint 传递给 Structured Streaming 引擎：

\begin{verbatim}
query = (
    enriched_stream_df.writeStream.foreachBatch(process_batch)
    .outputMode("append")
    .option("checkpointLocation", args.checkpoint_dir)
    .trigger(processingTime=f"{args.trigger_seconds} seconds")
    .start()
)
\end{verbatim}

其中 \texttt{checkpointLocation} 指向的目录用于保存每次微批次提交后的进度信息。在文件源模式下，Checkpoint 记录已处理的 JSON 文件名，确保同一文件不会被重复读取；在 Kafka 模式下，Checkpoint 记录每个 Partition 的消费偏移量（Offset），使系统重启后能够从上次中断的位置继续消费消息。

本项目中 Checkpoint 目录位于实时处理输出目录 \texttt{realtime/} 的子孙路径下。当系统使用 \texttt{--reset-output} 参数清空实时输出目录时，Checkpoint 也会随之一并被清除（通过 \texttt{safe\_reset\_output()} 中 \texttt{shutil.rmtree(realtime\_root)} 实现），从而确保每次全新的演示运行都从一致起点开始。这一设计与演示场景中的"全量重置"需求相匹配——在课程实验中，使用者通常希望每次运行得到独立、可复现的结果，而不是从上次中断位置继续处理。

此外，Checkpoint 机制与系统的安全目录清理策略形成了相互配合的关系。\texttt{safe\_reset\_output()} 在确认输出目录 basename 为 \texttt{"realtime"} 后，才执行删除操作。Checkpoint 目录作为输出目录的子目录也一并被安全清理。这种设计既保证了在正式运行环境中 Checkpoint 能够正常发挥容错作用，又保证了在实验演示场景中可以通过简单参数快速实现环境重置，较好地平衡了生产级可靠性和教学演示的便利性。

\subsection{数据源适配}

流处理引擎需要同时适配文件流模式和 Kafka 模式。两种模式对应不同的数据输入方式，但后续分析逻辑保持一致。为了实现这一目标，系统在数据接入层对不同来源的数据进行统一解析，使进入处理阶段的数据都具有相同字段结构。这样可以避免为文件源和 Kafka 源分别编写两套聚合统计和风险识别逻辑，提高系统的可维护性。

\subsubsection{文件源}

文件源是本系统默认采用的输入方式。在文件流模式下，事件生成器会按照批次将交易事件写入指定输入目录，每个批次对应一个 JSON Lines 文件。Spark Structured Streaming 通过 \texttt{readStream} 监听该目录，当目录中出现新的 JSON 文件时，系统会自动发现并读取新增内容。由于每一行 JSON 都表示一条完整交易事件，Spark 可以直接根据预定义 Schema 将其解析为结构化 DataFrame。

文件源模式的实现逻辑较为简洁，核心代码形式如下：

\begin{verbatim}
parsed_stream_df = (
    spark.readStream
    .schema(build_schema())
    .json(args.input_dir)
)
\end{verbatim}

其中，\texttt{build\_schema()} 用于构造事件字段 Schema，\texttt{args.input\_dir} 表示事件生成器写入 JSON 文件的目录。通过显式指定 Schema，系统可以避免 Spark 在流式读取过程中自动推断字段类型，从而提高解析稳定性。对于实时处理任务而言，显式 Schema 是一种更可靠的做法，因为流式数据可能分批到达，若依赖自动推断，字段类型可能因样本差异而出现不一致。

文件源模式适合教学演示和本地实验。用户可以直接查看输入目录中的批次文件，理解事件是如何被逐批写入的；也可以在调试时打开某个 JSON 文件，检查事件字段是否符合预期。对于本项目而言，文件源模式降低了实时处理链路的环境依赖，使系统能够在不启动 Kafka 服务的情况下完成完整的流式处理验证。

\subsubsection{Kafka 源}

Kafka 源用于支持更接近生产环境的实时数据处理模式。在 Kafka 模式下，事件生成器不再将事件写入本地文件，而是将每条交易事件序列化为 JSON 字符串，并发送到指定 Kafka Topic。流处理引擎通过 Spark Structured Streaming 的 Kafka Source 从该 Topic 中持续消费消息。

Kafka 消息中通常包含 \texttt{key}、\texttt{value}、\texttt{topic}、\texttt{partition}、\texttt{offset} 和 \texttt{timestamp} 等字段。本系统关注的是消息内容，因此需要将 Kafka 消息中的 \texttt{value} 字段取出，并转换为字符串。随后，系统使用 \texttt{from\_json()} 函数按照统一 Schema 将 JSON 字符串解析为结构化字段。

Kafka 源模式的核心处理逻辑如下：

\begin{verbatim}
raw_stream_df = (
    spark.readStream
    .format("kafka")
    .option("kafka.bootstrap.servers", args.bootstrap_servers)
    .option("subscribe", args.topic)
    .option("startingOffsets", args.starting_offsets)
    .option("failOnDataLoss", "false")
    .load()
)

parsed_stream_df = raw_stream_df.selectExpr(
    "CAST(value AS STRING) AS json_value"
)

parsed_stream_df = parsed_stream_df.select(
    from_json(col("json_value"), build_schema()).alias("event")
).select("event.*")
\end{verbatim}

在该过程中，\texttt{kafka.bootstrap.servers} 指定 Kafka Broker 地址，\texttt{subscribe} 指定需要消费的 Topic，\texttt{startingOffsets} 指定消费起始位置。系统设置 \texttt{failOnDataLoss} 为 \texttt{false}，可以避免 Kafka 日志因过期清理等原因导致流处理任务直接失败，从而提高演示和实验场景下的容错性。

Kafka 源模式的优势在于扩展性更强。对于真实生产系统而言，交易事件往往来自多个上游服务，例如订单系统、支付系统、用户行为日志系统等。Kafka 可以作为这些服务和下游分析系统之间的消息缓冲层，降低生产端和消费端之间的耦合度。本项目虽然默认使用文件流模式进行演示，但保留 Kafka 源适配能力，使系统具备向分布式实时处理架构迁移的基础。

\subsection{统一数据 Schema 设计}

为了保证文件模式和 Kafka 模式能够共用同一套处理逻辑，系统为实时交易事件定义了统一 Schema。无论事件来自 JSON 文件，还是来自 Kafka Topic 中的 JSON 消息，进入流处理引擎后都需要被解析为相同字段结构的 DataFrame。统一 Schema 是整个实时处理模块保持一致性的关键。

在流式处理场景中，Schema 的作用不仅是声明字段名称和数据类型，也起到约束输入数据格式的作用。如果没有明确 Schema，不同批次中字段类型的差异可能导致解析失败或统计结果异常。例如，\texttt{total\_spend} 必须被解析为数值类型，才能参与收入求和和平均消费计算；\texttt{discount\_applied} 必须被解析为布尔类型，才能统计折扣订单数量；\texttt{customer\_id} 和 \texttt{age} 等字段也需要保持整数类型，以保证后续处理的稳定性。

系统定义的实时事件 Schema 如表所示。

\begin{longtable}{lll}
\caption{实时事件 Schema 字段说明} \\
\toprule
字段名 & 数据类型 & 含义 \\
\midrule
\endfirsthead
\toprule
字段名 & 数据类型 & 含义 \\
\midrule
\endhead
\texttt{event\_id} & String & 事件唯一标识 \\
\texttt{event\_time} & String & 事件发生时间 \\
\texttt{customer\_id} & Integer & 客户编号 \\
\texttt{gender} & String & 性别 \\
\texttt{age} & Integer & 年龄 \\
\texttt{city} & String & 城市 \\
\texttt{membership\_type} & String & 会员等级 \\
\texttt{total\_spend} & Double & 消费金额 \\
\texttt{items\_purchased} & Integer & 购买件数 \\
\texttt{average\_rating} & Double & 平均评分 \\
\texttt{discount\_applied} & Boolean & 是否使用折扣 \\
\texttt{days\_since\_last\_purchase} & Integer & 距上次购买天数 \\
\texttt{satisfaction\_level} & String & 满意度等级 \\
\bottomrule
\end{longtable}

从字段设计看，Schema 同时包含事件追踪字段、客户属性字段和交易行为字段。\texttt{event\_id} 和 \texttt{event\_time} 用于标识和排序实时事件，使系统能够对批次内事件进行留存和追踪；\texttt{customer\_id}、\texttt{gender}、\texttt{age}、\texttt{city} 和 \texttt{membership\_type} 描述客户基础属性，为城市维度和会员维度分析提供依据；\texttt{total\_spend}、\texttt{items\_purchased}、\texttt{average\_rating}、\texttt{discount\_applied} 和 \texttt{days\_since\_last\_purchase} 描述交易与行为特征，是 KPI 计算和风险识别的主要输入；\texttt{satisfaction\_level} 则表示客户满意度状态，用于辅助判断客户是否存在流失风险。

统一 Schema 还增强了系统的可扩展性。如果后续需要增加新的分析字段，例如商品类别、支付方式、订单渠道或地区编码，只需要在事件生成器和 Schema 定义中同步扩展字段，再在流处理逻辑中增加相应统计即可。由于文件源和 Kafka 源都依赖同一套 Schema，字段扩展不需要重复修改两套处理逻辑。

\subsection{数据增强逻辑}

流处理引擎在读取原始事件后，会对数据进行增强处理。数据增强的目的是在原始字段基础上构造更具业务含义的派生字段，使后续聚合统计和风险识别更加直接。本系统主要新增两个派生字段：风险客户标记 \texttt{is\_at\_risk} 和高价值订单标记 \texttt{high\_value\_order}。

数据增强逻辑由 \texttt{enrich\_stream()} 函数实现。该函数接收解析后的事件流 DataFrame，并通过 Spark SQL 表达式新增派生列。由于这些派生列是在流处理阶段实时计算得到的，因此能够随着每个微批次的新事件动态更新。相比在事件生成阶段直接写死这些标记，在流处理阶段计算派生字段能够使业务规则更加集中，也便于后续调整识别逻辑。

\subsubsection{风险客户标记}

风险客户标记用于识别可能存在流失风险或需要运营关注的客户。在电商业务中，客户是否存在风险通常不能只通过单一指标判断。长时间未购买可能意味着客户活跃度下降，评分较低或满意度较低则可能意味着客户体验不佳。系统将这两类因素结合起来，构造风险客户识别规则。

风险客户的判定逻辑可以形式化表示为：

\[
\begin{aligned}
\text{is\_at\_risk} =\; &(\text{days\_since\_last\_purchase} > 30) \\
\land\; &(\text{satisfaction\_level} = \text{Unsatisfied} \\
&\lor\; \text{average\_rating} < 3.8)
\end{aligned}
\]

该规则包含两个层面的约束。时间维度要求 \texttt{days\_since\_last\_purchase} 大于 30，表示客户已经较长时间没有购买行为。情感维度要求客户满意度为 \texttt{Unsatisfied}，或者平均评分低于 3.8，表示客户对平台或商品体验可能不佳。只有当客户同时满足较长时间未购买以及满意度或评分较低时，系统才将其标记为风险客户。

这一规则具有较强的业务解释性。若客户只是短期未购买，但评分较高且满意度较好，则不应过早判定为风险客户；若客户评分偏低但近期仍有购买行为，也可能只是一次短期体验波动。将时间因素和满意度因素结合，可以减少单一指标带来的误判，使风险客户识别更加稳健。

在代码层面，系统通过 Spark 的 \texttt{when}、\texttt{col} 和 \texttt{lit} 等函数构造布尔表达式，并将判断结果保存为 \texttt{is\_at\_risk} 字段。后续风险客户明细输出会筛选该字段为 \texttt{True} 的事件，并按照距上次购买天数和平均评分进行排序，使风险程度较高的客户优先展示。

\begin{figure}[H]
\centering
\includegraphics[width=0.65\textwidth]{figures/pic4.png}
\caption{增强型风险客户识别规则决策流程：在传统"距上次购买 > 30 天且满意度 = Unsatisfied"的单条件基础上，引入"评分 < 3.8"作为替代判定条件，实现时间维度与情感维度的双重覆盖，提升风险客户识别的灵敏度与覆盖面}
\label{fig:risk_detection}
\end{figure}

\subsubsection{高价值订单标记}

高价值订单标记用于识别单笔消费金额较高的交易事件。在电商运营中，高价值订单往往代表较强的购买能力和较高的业务价值。统计不同批次或不同会员类型中的高价值订单数量，可以帮助分析客户消费层级和会员价值差异。

本系统将单笔消费金额大于或等于 1000 的交易标记为高价值订单。该规则可以形式化表示为：

\[
high\_value\_order =
\begin{cases}
1, & total\_spend \geq 1000 \\
0, & total\_spend < 1000
\end{cases}
\]

这里将 \texttt{high\_value\_order} 设计为数值型标记，而不是布尔型标记，主要是为了方便后续聚合统计。在城市维度或会员维度分析中，系统可以直接对该字段求和，得到高价值订单数量。例如，在会员维度聚合中，对 Gold、Silver 和 Bronze 等会员类型分别统计 \texttt{high\_value\_order} 的总和，即可比较不同会员群体贡献高价值订单的能力。

高价值订单标记虽然规则简单，但在实时 KPI 中具有明确意义。随着新批次事件不断到达，系统可以持续统计当前批次高价值订单数量，辅助判断最新交易质量。如果某一批次事件数量较多但高价值订单数量较少，可能说明该批次以低客单价交易为主；如果高价值订单数量明显增加，则可能反映当前批次中出现了更高价值的消费行为。

\subsection{微批次处理流程}

本系统通过 \texttt{foreachBatch} 机制实现微批次处理。\texttt{foreachBatch} 是 Spark Structured Streaming 中用于自定义批次处理逻辑的重要接口，它允许开发者在每个微批次到达时，将该批次数据作为一个静态 DataFrame 传入自定义函数。这样，系统可以在每个批次内部使用常规 DataFrame 操作完成排序、聚合、筛选、转换和写出等复杂逻辑。

在本项目中，微批次处理逻辑集中封装在 \texttt{process\_batch(batch\_df, batch\_id)} 函数中。其中，\texttt{batch\_df} 表示当前微批次的数据，\texttt{batch\_id} 表示当前批次编号。每当 Spark Structured Streaming 检测到新事件并形成微批次后，都会调用该函数执行完整分析流程。该设计使实时处理逻辑集中在一个入口中，便于维护和调试。

\subsubsection{空批次过滤}

流式处理任务在运行过程中可能会遇到空批次。空批次通常表示当前触发周期内没有新的输入事件。如果对空批次继续执行后续聚合和写文件操作，可能会生成空结果文件，甚至影响仪表盘展示效果。因此，系统在每个批次开始时都会检查事件数量。

其基本逻辑如下：

\begin{verbatim}
event_count = batch_df.count()
if event_count == 0:
    return
\end{verbatim}

当 \texttt{event\_count} 为 0 时，系统直接结束当前批次处理，不写入任何结果文件。该机制能够避免无效输出，也使系统在没有新数据到达时保持稳定运行。对于实时处理任务而言，空批次过滤是一种必要的健壮性设计。

\subsubsection{事件留存}

对于非空批次，系统会对事件进行留存。事件留存的目标是保存原始或增强后的交易明细，使后续能够追踪每个批次中具体包含哪些事件。本系统采用三类保存策略：按批次保存、最新批次快照和全量事件追加。

按批次保存是指将每个微批次中的事件写入独立 CSV 文件，例如 \texttt{events/batch\_000.csv}、\texttt{events/batch\_001.csv} 等。该方式便于逐批追溯，当某个批次的统计结果异常时，可以直接查看对应批次的事件明细，分析异常来源。

最新批次快照是指将当前批次事件覆盖写入 \texttt{latest\_events.csv}。该文件始终保存最近一次处理批次的完整事件列表，主要供 HTML 仪表盘展示使用。由于仪表盘通常只需要展示最新状态，因此采用覆盖写入方式可以避免读取过多历史数据。

全量事件追加是指将每个批次事件追加写入 \texttt{stream\_events.csv}。该文件记录整个实时处理会话中的累计事件历史，能够用于后续复盘、统计和验证。追加写入方式可以保留完整处理过程，使系统不仅具备实时展示能力，也具备历史追踪能力。

三类留存策略分别面向不同需求：批次归档面向调试追溯，最新快照面向实时展示，全量追加面向历史记录。通过将三者结合，系统能够兼顾实时性和可追溯性。

\subsubsection{城市维度聚合}

城市维度聚合用于分析不同城市在当前批次中的交易表现。电商客户来自不同城市，不同城市可能在消费能力、购买频率和订单规模方面存在差异。通过城市维度统计，系统可以观察最新批次中各城市的事件数量、平均消费金额和购买商品总件数。

城市维度聚合的主要指标包括：
\begin{itemize}
    \item \texttt{event\_count}：各城市在当前批次中的交易事件数量；
    \item \texttt{avg\_total\_spend}：各城市当前批次的平均消费金额；
    \item \texttt{total\_items\_purchased}：各城市当前批次购买商品总件数。
\end{itemize}

其处理逻辑可以概括为按 \texttt{city} 分组，并对消费金额和购买数量进行聚合计算：

\begin{verbatim}
city_summary_pdf = (
    ordered_df.groupBy("city")
    .agg(
        count("*").alias("event_count"),
        spark_round(avg("total_spend"), 2).alias("avg_total_spend"),
        spark_sum("items_purchased").alias("total_items_purchased"),
    )
    .orderBy("city")
    .toPandas()
)
\end{verbatim}

城市聚合结果会被写入 \texttt{city\_summary\_latest.csv}。由于该文件表示最新批次的城市维度快照，因此采用覆盖写入方式。HTML 仪表盘可以读取该文件，并将各城市平均消费金额展示为柱状图，使用户直观比较不同城市的消费表现。

\subsubsection{会员维度聚合}

会员维度聚合用于分析不同会员等级客户在当前批次中的行为差异。电商平台通常会根据客户价值、消费频率或权益等级划分会员类型，例如 Gold、Silver 和 Bronze。不同会员等级在消费金额、评分水平和高价值订单数量方面可能存在明显差异，因此会员维度分析对于理解用户分层具有重要意义。

会员维度聚合的主要指标包括：
\begin{itemize}
    \item \texttt{event\_count}：各会员等级在当前批次中的交易事件数量；
    \item \texttt{avg\_total\_spend}：各会员等级当前批次平均消费金额；
    \item \texttt{avg\_rating}：各会员等级当前批次平均评分；
    \item \texttt{high\_value\_orders}：各会员等级当前批次高价值订单数量。
\end{itemize}

其处理逻辑可以概括为按 \texttt{membership\_type} 分组，并分别计算事件数量、平均消费、平均评分和高价值订单总数：

\begin{verbatim}
membership_summary_pdf = (
    ordered_df.groupBy("membership_type")
    .agg(
        count("*").alias("event_count"),
        spark_round(avg("total_spend"), 2).alias("avg_total_spend"),
        spark_round(avg("average_rating"), 2).alias("avg_rating"),
        spark_sum("high_value_order").alias("high_value_orders"),
    )
    .orderBy("membership_type")
    .toPandas()
)
\end{verbatim}

会员聚合结果会写入 \texttt{membership\_summary\_latest.csv}，用于仪表盘中的会员快照展示。通过该文件，用户可以观察当前批次中不同会员类型的交易活跃度、消费能力和订单价值。例如，如果 Gold 会员的平均消费和高价值订单数量较高，则说明高等级会员对当前批次收入贡献较大；如果某类会员平均评分偏低，则可能需要进一步分析其体验问题。

\subsubsection{风险客户识别}

风险客户识别是本系统实时处理模块的重要业务功能。系统会筛选出 \texttt{is\_at\_risk} 字段为 \texttt{True} 的事件，并生成风险客户明细表。该表用于展示当前批次中可能存在流失风险或需要运营关注的客户。

风险客户明细包含事件时间、客户编号、城市、会员类型、消费金额、平均评分、距上次购买天数和满意度等级等字段。这些字段能够帮助运营人员理解客户被判定为风险客户的原因。例如，某客户距上次购买天数较长且评分偏低，说明其近期活跃度和体验评价均不理想；某客户属于高等级会员但被标记为风险客户，则可能具有更高的挽回价值。

系统在输出风险客户明细时，会按照风险程度进行排序。排序规则主要包括：
\begin{itemize}
    \item 距上次购买天数越长，客户风险越高；
    \item 平均评分越低，客户风险越高。
\end{itemize}

对应逻辑可以概括为：

\begin{verbatim}
risk_pdf = (
    ordered_df.filter(col("is_at_risk"))
    .select(
        "event_time", "customer_id", "city", "membership_type",
        "total_spend", "average_rating",
        "days_since_last_purchase", "satisfaction_level",
    )
    .orderBy(
        col("days_since_last_purchase").desc(),
        col("average_rating").asc()
    )
    .toPandas()
)
\end{verbatim}

排序后的风险客户结果会写入 \texttt{at\_risk\_customers\_latest.csv}。HTML 仪表盘读取该文件后，可以展示风险客户 Top 列表，使用户在查看实时报告时能够快速定位需要关注的客户群体。

\subsubsection{KPI 指标汇总}

KPI 指标汇总用于概括每个微批次的整体业务表现。与城市聚合和会员聚合侧重局部维度不同，批次 KPI 更关注当前批次的整体交易情况和系统处理进度。每处理一个非空批次，系统都会生成一条 KPI 记录，并追加写入 \texttt{kpi\_history.csv}，同时覆盖写入 \texttt{latest\_batch\_metrics.csv}。

每条 KPI 记录主要包含以下指标：
\begin{itemize}
    \item \texttt{batch\_id}：当前批次编号；
    \item \texttt{processed\_ batches}：已经处理的批次数量；
    \item \texttt{event\_count}：当前批次事件数量；
    \item \texttt{cumulative\_event\_count}：累计事件数量；
    \item \texttt{total\_revenue}：当前批次总收入；
    \item \texttt{avg\_spend}：当前批次平均消费金额；
    \item \texttt{avg\_rating}：当前批次平均评分；
    \item \texttt{discounted\_orders}：当前批次折扣订单数量；
    \item \texttt{high\_value\_orders}：当前批次高价值订单数量；
    \item \texttt{at\_risk\_count}：当前批次风险客户数量。
\end{itemize}

KPI 计算主要基于当前批次事件明细完成。例如，总收入由 \texttt{total\_spend} 求和得到，平均消费由消费金额均值表示，折扣订单数量由 \texttt{discount\_applied} 为真的事件数量得到，高价值订单数量由 \texttt{high\_value\_order} 求和得到，风险客户数量则由风险客户明细行数得到。

KPI 指标汇总具有两方面作用。对于系统运行而言，KPI 历史可以反映处理进度和事件累计情况，便于判断流处理任务是否正常执行；对于业务分析而言，KPI 指标可以快速概括最新交易状态，例如当前批次收入是否较高、平均评分是否下降、风险客户数量是否增加等。由于 \texttt{kpi\_history.csv} 采用追加写入方式，系统可以保留多个批次的时间序列记录，为后续趋势分析提供基础。

\begin{figure}[H]
\centering
\includegraphics[width=0.60\textwidth]{figures/pic2.png}
\caption{\texttt{process\_batch()} 微批次处理内部流程图：从空批次过滤开始，依次执行事件留存、城市维度聚合、会员维度聚合、风险客户识别、KPI 指标汇总和仪表盘渲染，最后根据停止条件判断是否继续等待下一微批次}
\label{fig:microbatch_flow}
\end{figure}

\subsection{流控制与停止策略}

在真实生产环境中，流处理任务通常会长期运行，持续等待和处理新事件。然而，在课程实验和本地演示场景中，系统不适合无限运行，否则会增加资源占用和操作复杂度。因此，本项目为流处理引擎设计了明确的流控制与停止策略，使任务能够在完成指定处理目标后自动结束。

系统主要采用两类停止条件：批次上限和空闲超时。批次上限表示当系统已经处理指定数量的非空批次后，自动停止流处理任务。该机制适合实验演示，因为用户可以提前设置需要观察的批次数，例如处理 5 个批次后结束。这样既能完整展示流式处理过程，也能避免任务持续占用资源。

空闲超时表示当系统在一段时间内没有接收到新事件时，自动停止流处理任务。该机制可以避免事件生成器已经结束但流处理引擎仍然持续等待的情况。例如，当所有批次事件都已写入输入目录并处理完毕后，如果后续没有新文件产生，系统会在空闲时间达到阈值后结束运行。空闲超时机制提高了实验流程的自动化程度，使用户无需手动终止流处理进程。

为了跟踪处理状态，系统使用状态变量记录已处理批次数、累计事件数量、是否需要停止以及最近一次批次完成时间等信息。其状态结构可以概括为：

\begin{verbatim}
state = {
    "processed_batches": 0,
    "cumulative_event_count": 0,
    "should_stop": False,
    "last_batch_completed_at": None,
}
\end{verbatim}

其中各字段的作用如下：
\begin{itemize}
    \item \texttt{processed\_batches}：判断是否达到批次上限；
    \item \texttt{cumulative\_event\_count}：记录累计处理事件数；
    \item \texttt{should\_stop}：标识是否需要结束查询；
    \item \texttt{last\_batch\_completed\_at}：配合空闲超时机制，判断是否长时间无新数据到达。
\end{itemize}

流控制与停止策略使本系统更适合课程项目运行环境。用户可以启动事件生成器和流处理引擎，等待系统自动完成若干批次处理并停止，然后直接查看输出结果。相比需要人工观察并手动终止的方式，该策略减少了操作步骤，也降低了因忘记关闭任务而造成资源浪费的可能性。

\subsection{安全机制设计}

实时处理引擎在运行过程中会生成大量输出文件，包括批次事件文件、最新事件快照、KPI 历史、城市聚合结果、会员聚合结果、风险客户列表和仪表盘文件等。为了便于重复运行实验，系统支持在启动时清空旧的实时处理输出目录。然而，自动删除目录是一项具有风险的操作，如果路径设置错误，可能会误删项目中的其他重要文件。因此，系统在清空输出目录时加入了路径安全检查机制。

安全机制的核心思想是：系统只允许删除明确命名为 \texttt{realtime} 的输出目录，而不会对任意路径执行删除操作。其逻辑可以概括为：

\begin{verbatim}
def safe_reset_output(output_dir, checkpoint_dir):
    realtime_root = os.path.abspath(output_dir)
    if os.path.basename(realtime_root) != "realtime":
        raise ValueError(
            f"Refusing to delete unexpected output directory: {output_dir}"
        )
    if os.path.exists(realtime_root):
        shutil.rmtree(realtime_root)
\end{verbatim}

该函数先将输出目录转换为绝对路径，再检查路径最后一级目录名是否为 \texttt{realtime}。如果目录名不符合预期，系统会拒绝删除并抛出异常。只有在确认目标目录确实是实时处理输出目录后，系统才会执行删除操作。通过这种方式，即使用户误传了项目根目录或其他重要路径，程序也不会直接清空该目录。

安全机制设计体现了系统工程实现中的防御性编程思想。对于数据处理系统而言，路径、文件和目录操作非常常见，但也容易引发误删、覆盖和数据丢失等问题。在本项目中，虽然实时处理输出主要用于实验和展示，但加入路径安全检查仍然能够显著提高系统健壮性和可靠性。

除路径安全检查外，系统在其他环节也体现了稳健性设计。例如，空批次过滤避免生成无意义文件；统一 Schema 减少字段解析错误；文件追加函数避免重复写入表头；批次上限和空闲超时防止流处理任务长时间挂起。这些机制共同保证了流处理引擎在实验环境中的稳定运行，也使系统更符合可重复执行和可维护的工程要求。

\section{报表生成与可视化设计}

\subsection{输出目录结构}

报表生成与可视化模块位于实时处理链路的末端，主要负责对流处理引擎生成的结构化结果进行整理和展示。实时处理过程中会产生多类输出文件，包括事件明细、批次 KPI、城市维度聚合、会员维度聚合、风险客户列表以及可视化报告等。为了便于管理和查看，系统将所有实时处理结果统一保存到 \texttt{Reports/realtime/} 目录下，并在该目录中进一步组织不同类型的输出内容。

从目录结构上看，\texttt{Reports/realtime/} 既是流处理结果的存储位置，也是报表生成器读取数据的基础目录。流处理引擎在每个微批次完成后，会将当前批次的统计结果写入该目录下的多个 CSV 文件；报表生成器随后读取这些 CSV 文件，生成 HTML 仪表盘和 Markdown 摘要。通过这种方式，计算模块和展示模块之间通过文件系统进行解耦，流处理引擎只需要负责输出结构化结果，报表生成器只需要负责读取结果并渲染页面，两者职责清晰。

在 \texttt{Reports/realtime/} 目录下，系统还设置了 \texttt{events/} 子目录，用于保存每个微批次的事件归档文件。每个批次对应一个独立 CSV 文件，例如 \texttt{batch\_000.csv}、\texttt{batch\_001.csv} 等。该设计使用户能够追溯任意批次的输入事件明细。当某个批次的 KPI 指标或聚合结果出现异常时，可以直接打开对应批次文件，检查该批次中具体包含哪些交易事件，从而定位问题来源。

除批次事件归档外，实时处理目录还保存若干快照型和历史型 CSV 文件。快照型文件通常采用覆盖写入方式，例如 \texttt{latest\_events.csv}、\texttt{latest\_batch\_metrics.csv}、\texttt{city\_summary\_latest.csv} 和 \texttt{membership\_summary\_latest.csv} 等。这类文件始终表示最新批次或最新状态，适合被 HTML 仪表盘直接读取。历史型文件则采用追加写入方式，例如 \texttt{kpi\_history.csv} 和 \texttt{stream\_events.csv}，用于保存整个流式处理过程中的累计记录，便于后续复盘和趋势分析。

整体而言，输出目录结构兼顾了实时展示和历史追踪两类需求。快照文件保证仪表盘能够快速读取最新结果，历史文件保证处理过程能够被完整保留，批次归档文件则为逐批排查和结果验证提供依据。该设计使实时处理模块不仅能够”实时展示结果”，也能够”回溯分析过程”。

\begin{figure}[H]
\centering
\includegraphics[width=0.85\textwidth]{figures/realtime_output_dir.png}
\caption{实时处理输出目录结构（\texttt{Reports/realtime/}），包含 7 类 CSV 文件、HTML 仪表盘、Markdown 摘要、批次事件归档子目录和 Checkpoint 目录}
\label{fig:output_dir}
\end{figure}

\subsection{CSV 输出文件设计}

CSV 文件是本系统实时处理结果的主要结构化载体。相比直接将结果写入 HTML 页面，CSV 文件具有更好的通用性和可复用性。它既可以被报表生成器读取，用于生成仪表盘和摘要，也可以被用户直接打开检查，或被其他数据分析工具继续处理。因此，本系统将流处理阶段产生的主要分析结果都以 CSV 文件形式保存。

\begin{table}[H]
\centering
\caption{实时处理输出文件}
\begin{tabular}{lll}
\toprule
输出文件 & 写入策略 & 用途 \\
\midrule
\texttt{kpi\_history.csv} & 追加 & 批次级 KPI 时间序列 \\
\texttt{latest\_batch\_metrics.csv} & 覆盖 & 最新批次 KPI 快照 \\
\texttt{latest\_events.csv} & 覆盖 & 最新批次事件明细 \\
\texttt{stream\_events.csv} & 追加 & 全量事件历史 \\
\texttt{city\_summary\_latest.csv} & 覆盖 & 城市维度聚合结果 \\
\texttt{membership\_summary\_latest.csv} & 覆盖 & 会员维度聚合结果 \\
\texttt{at\_risk\_customers\_latest.csv} & 覆盖 & 风险客户明细 \\
\texttt{events/batch\_XXX.csv} & 按批写入 & 单批次事件归档 \\
\bottomrule
\end{tabular}
\end{table}

\texttt{kpi\_history.csv} 用于保存批次级 KPI 时间序列。系统每处理一个非空微批次，就会向该文件追加一条新的 KPI 记录。该文件中的每一行对应一个批次，记录内容包括批次编号、处理时间、当前批次事件数、累计事件数、总收入、平均消费、平均评分、折扣订单数、高价值订单数和风险客户数量等。由于该文件采用追加写入方式，因此能够完整反映实时处理任务从开始到结束的运行过程。后续若需要分析不同批次之间 KPI 的变化趋势，可以直接基于该文件进行计算。

\texttt{latest\_batch\_metrics.csv} 用于保存最新批次的 KPI 快照。与 \texttt{kpi\_history.csv} 不同，该文件每次都会被最新批次结果覆盖，因此始终只包含最近一次处理批次的核心指标。HTML 仪表盘中的摘要卡片区主要依赖该文件展示当前最新状态，例如最新批次收入、处理批次数、累计事件数和风险客户数量等。采用覆盖写入方式可以减少仪表盘读取数据时的处理成本，使展示逻辑更加简单。

\texttt{latest\_events.csv} 用于保存最新批次的事件明细。该文件包含当前微批次中所有交易事件及其增强后的字段信息，例如客户编号、城市、会员类型、消费金额、评分、是否为风险客户以及是否为高价值订单等。该文件主要用于仪表盘中的“最新事件明细表”，帮助用户查看系统最近处理了哪些交易事件。由于实时展示通常更关注最新批次状态，因此该文件采用覆盖写入方式。

\texttt{stream\_events.csv} 用于保存整个实时处理过程中的全量事件历史。系统每处理一个批次，就会将该批次中的事件追加到该文件中。与 \texttt{latest\_events.csv} 只保留最新批次不同，\texttt{stream\_events.csv} 可以完整记录所有已处理事件，因此适合用于后续复盘、累计统计和结果验证。例如，可以通过该文件检查累计事件数是否与各批次事件数之和一致，也可以重新计算全局收入和评分分布。

\texttt{city\_summary\_latest.csv} 用于保存城市维度的最新聚合结果。该文件按城市分组，统计当前批次中各城市的交易事件数量、平均消费金额和购买商品总件数。城市维度聚合能够反映不同地区客户在最新批次中的消费表现，也为仪表盘中的城市平均消费柱状图提供数据来源。由于该文件表示最新批次快照，因此采用覆盖写入方式。

\texttt{membership\_summary\_latest.csv} 用于保存会员维度的最新聚合结果。该文件按会员类型分组，统计不同会员等级的事件数量、平均消费金额、平均评分和高价值订单数量。会员维度结果有助于比较 Gold、Silver、Bronze 等不同客户群体的实时表现。例如，高等级会员是否贡献了更多高价值订单，某一类会员的平均评分是否偏低，都可以通过该文件进行观察。

\texttt{at\_risk\_customers\_latest.csv} 用于保存最新批次识别出的风险客户明细。该文件筛选出满足风险判断条件的客户事件，并保留客户编号、城市、会员类型、消费金额、评分、距上次购买天数和满意度等级等关键字段。由于风险客户识别具有较强业务意义，该文件可以帮助运营人员快速定位需要关注的客户。仪表盘中的风险客户 Top 列表也主要基于该文件生成。

\texttt{events/batch\_XXX.csv} 用于保存单个批次的事件归档。每个批次都会生成一个独立文件，文件名中包含批次编号。该设计使实时处理过程具有较强的可追踪性。当需要验证某一批次的 KPI 是否计算正确时，可以直接查看对应批次文件，并与聚合结果进行比对。相比只保存最新快照，批次归档能够更完整地保留流处理过程中的中间数据。

从写入策略看，本系统有意识地区分了“覆盖写入”和“追加写入”。覆盖写入适用于最新状态类文件，可以保证仪表盘读取到的是当前最新结果；追加写入适用于历史记录类文件，可以保留完整处理过程。这种设计避免了所有结果都不断累积导致仪表盘读取成本过高的问题，也避免了只保存最新结果导致历史信息丢失的问题。

\subsection{HTML 仪表盘设计}

HTML 仪表盘是本系统实时分析结果的主要可视化展示形式。相比单纯查看 CSV 文件，仪表盘能够更加直观地呈现关键指标、聚合结果和风险客户信息。用户只需要在浏览器中打开 \texttt{dashboard.html}，即可查看当前实时处理任务的整体运行状态和最新分析结果。

仪表盘采用静态 HTML 文件形式生成，不依赖额外前端框架或服务器环境。该设计适合课程项目和本地实验场景，因为用户无需配置 Web 服务，也无需安装复杂的前端依赖。报表生成器会在每个微批次处理完成后重新渲染 HTML 文件，使页面内容随着最新 CSV 结果更新。虽然该方式不是严格意义上的前端实时推送，但对于微批次流处理演示而言，已经能够清晰展示结果随批次刷新的过程。

在页面结构上，HTML 仪表盘主要包含摘要卡片区、图表展示区和数据表格区。摘要卡片区位于页面顶部，用于展示最核心的实时指标，包括已处理批次数、累计事件数、最新批次收入和风险客户数量。这些指标能够帮助用户快速判断系统当前运行状态。例如，累计事件数可以反映流处理进度，最新批次收入可以反映当前交易规模，风险客户数量则可以提示当前批次是否存在需要关注的客户群体。

城市平均消费柱状图用于展示不同城市在最新批次中的平均消费金额。报表生成器读取 \texttt{city\_summary\_latest.csv} 后，将各城市平均消费金额转换为 HTML 中的柱状图形式。通过柱状图，用户可以直观看到哪些城市的平均消费较高，哪些城市的消费表现相对较低。相比表格形式，柱状图更适合进行横向比较，也更适合在报告展示中说明区域消费差异。

会员类型事件分布图用于展示不同会员等级在最新批次中的事件数量。该部分基于 \texttt{membership\_summary\_latest.csv} 生成，能够反映当前批次中各会员类型的活跃程度。例如，如果某一批次中 Gold 会员事件数量较多，说明高等级会员在该批次中较为活跃；如果 Bronze 会员占比较高，则可能说明当前交易更多来自普通客户群体。该图表与会员快照表结合，可以进一步分析会员等级与消费金额、评分和高价值订单之间的关系。

批次 KPI 历史表用于展示最近若干个批次的关键指标变化。该表基于 \texttt{kpi\_history.csv} 生成，通常会按照批次时间倒序展示最近记录。通过 KPI 历史表，用户可以观察不同批次之间的事件数、收入、平均消费、平均评分和风险客户数量变化。该部分使仪表盘不仅能够展示最新状态，也能够提供一定的时间序列视角。

风险客户 Top 列表是仪表盘中业务价值较高的部分。该列表基于 \texttt{at\_risk\_customers\_latest.csv} 生成，展示当前批次中风险程度较高的客户。系统在流处理阶段已经按照距上次购买天数降序和平均评分升序进行排序，因此列表前面的客户通常具有更高关注优先级。通过该表，用户可以快速查看风险客户的城市、会员类型、消费金额、评分和满意度等级，为后续客户维护或营销干预提供依据。

最新事件明细表用于展示当前批次中的交易事件。该表基于 \texttt{latest\_events.csv} 生成，可以帮助用户了解最近一次微批次具体处理了哪些数据。对于系统验证而言，最新事件明细表非常重要，因为它能够将输入事件与输出指标联系起来。例如，当摘要卡片显示某批次收入较高时，用户可以查看最新事件明细，判断是否存在多个高消费订单。

会员快照表用于展示不同会员类型在最新批次中的聚合指标，包括事件数量、平均消费、平均评分和高价值订单数量等。该表与会员类型事件分布图共同构成会员维度分析部分。图表用于展示事件数量差异，表格则提供更完整的数值信息。通过会员快照表，用户可以更细致地比较不同会员等级的消费能力和满意度水平。

在视觉设计方面，仪表盘采用内联 CSS 实现布局和样式，避免依赖外部样式文件。页面通常使用卡片式布局和网格结构，将核心指标、图表和表格分区展示。这样既保证页面结构清晰，也便于在不同屏幕尺寸下查看。由于所有样式都写入单个 HTML 文件，用户可以直接复制或打开该文件，不需要额外资源支持。

总体来看，HTML 仪表盘承担了将结构化数据转化为可读报告的作用。它将实时处理产生的多个 CSV 文件整合到同一页面中，使用户能够从整体指标、城市维度、会员维度、风险客户和事件明细等多个角度理解当前实时处理结果。对于课程项目而言，该仪表盘也增强了系统展示效果，使实时处理链路的运行结果更加直观。

\begin{figure}[H]
\centering
\includegraphics[width=\textwidth]{figures/dashboard_full.png}
\caption{HTML 实时仪表盘全貌（\texttt{dashboard.html}），包含摘要卡片（已处理批次数、累计事件数、最新收入、风险客户数）、城市平均消费柱状图、会员类型事件分布、批次 KPI 历史表、风险客户 Top 列表、最新事件明细表和会员快照表}
\label{fig:dashboard_full}
\end{figure}

\subsection{Markdown 摘要设计}

除 HTML 仪表盘外，系统还会生成 Markdown 摘要文件 \texttt{summary.md}。与 HTML 仪表盘相比，Markdown 摘要更加简洁，适合快速查看实时处理任务的运行概况，也便于在实验记录或报告中引用。该文件通常包含报告生成时间、累计处理批次数、累计事件数、最新批次收入、平均消费、风险客户数量以及主要输出文件说明等内容。

Markdown 摘要的生成依赖最新批次 KPI 文件和相关输出结果。报表生成器读取最新指标后，将关键数值组织为文本段落或列表形式。由于 Markdown 文件本身是纯文本格式，因此具有良好的可读性和可移植性。用户可以在代码编辑器、Markdown 阅读器或版本管理系统中直接查看该文件，不需要打开浏览器或加载额外样式。

报告生成时间用于标识摘要文件的刷新时刻。实时处理任务会在每个微批次完成后更新摘要，因此生成时间可以帮助用户判断当前摘要是否对应最新结果。若报告生成时间明显早于实验结束时间，说明摘要可能尚未刷新或处理任务未正常运行。

累计批次数和累计事件数用于反映实时处理任务的整体进度。累计批次数表示系统已经成功处理了多少个非空微批次，累计事件数表示所有已处理批次中事件数量的总和。这两个指标有助于验证事件生成器和流处理引擎是否协调运行。例如，如果事件生成器发布了多个批次，而摘要中累计批次数偏少，则可能说明流处理任务没有读取到全部输入数据。

最新批次收入和平均消费用于概括当前最新交易批次的业务表现。最新批次收入反映当前批次交易金额总和，平均消费则反映单条事件的平均消费水平。二者结合可以帮助用户快速判断当前批次是以高客单价交易为主，还是以较多低额交易为主。

风险客户数量用于提示当前批次中需要关注的客户规模。该指标来源于风险客户识别结果，能够帮助用户快速判断最新事件中是否存在较明显的客户流失风险。如果风险客户数量较高，用户可以进一步打开 HTML 仪表盘或 \texttt{at\_risk\_customers\_latest.csv} 查看具体客户明细。

输出文件说明是 Markdown 摘要中的辅助信息，用于列出实时处理模块生成的主要文件及其作用。通过该部分，用户可以快速了解每个 CSV 文件、HTML 文件和 Markdown 文件的用途。例如，\texttt{kpi\_history.csv} 用于查看批次 KPI 时间序列，\texttt{dashboard.html} 用于查看可视化仪表盘，\texttt{stream\_events.csv} 用于查看全量事件历史。该设计降低了用户理解输出目录结构的成本。

Markdown 摘要的价值在于提供轻量级运行概览。HTML 仪表盘适合展示较丰富的图表和表格，而 Markdown 摘要适合快速阅读和结果归档。二者相互补充，使系统的输出既有较强的可视化效果，也有便于记录和引用的文本说明。对于期末报告而言，Markdown 摘要中的核心指标还可以作为实验结果验证部分的数据来源之一。

\section{Kafka 模式扩展设计}

\subsection{原始 Kafka 链路}

在实时数据处理系统中，Kafka 常被用于连接数据生产端和数据消费端。它可以作为中间消息队列，将上游持续产生的事件数据传递给下游流式计算模块。对于电商客户行为分析场景而言，用户交易、评价、优惠使用和会员行为等事件都可以被视为持续产生的消息流。通过 Kafka 构建生产者与消费者链路，可以使数据生成模块和数据处理模块之间实现解耦，从而提高系统的扩展性和稳定性。

本项目在设计实时处理模块时，参考了原有的 Kafka Producer 和 Kafka Consumer 实现。原始 Kafka Producer 的主要功能是生成模拟客户行为数据，并将其发送到指定 Kafka Topic。该 Producer 会构造客户编号、性别、年龄、城市、会员类型、消费金额、购买数量、评分、折扣使用情况和满意度等级等字段，并按照一定时间间隔将数据持续写入 Kafka。通过这种方式，Producer 模拟了电商平台中交易事件不断产生并进入消息队列的过程。

原始 Kafka Consumer 则负责从 Kafka Topic 中消费消息，并使用 Spark Structured Streaming 对消息进行实时处理。Consumer 启动后会持续监听指定 Topic，当 Kafka 中出现新消息时，Spark 会将这些消息读入流式 DataFrame，并根据预设解析逻辑提取字段。随后，Consumer 可以对客户行为数据进行实时聚合，例如按照会员类型、年龄组和城市等维度统计客户数量、消费情况或其他指标。通过 Producer 与 Consumer 的配合，原始链路已经具备了基本的实时数据生成、消息传递和流式分析能力。

不过，原始 Kafka 链路在数据生成方式、消息格式和处理能力上存在可改进空间。以下从代码层面对原始链路与本项目适配版本进行具体对比。

\subsubsection{Producer 端对比}

原始 Producer（\texttt{kafka\_producer\_ecommerce.py}）采用纯随机生成方式：

\begin{verbatim}
customer_id = random.randint(101, 999)
city = random.choice(['New York', 'Los Angeles', ...])
total_spend = round(random.uniform(400, 1500), 2)
average_rating = round(random.uniform(3.0, 4.7), 1)
message = f"{customer_id},{gender},...,{satisfaction_level}"
\end{verbatim}

该方式存在三个主要局限：(1) 各字段独立随机生成，客户编号、城市、会员类型与消费行为之间缺乏业务关联，例如可能出现低消费 Gold 会员或评分与满意度不匹配等不合理组合；(2) 消息采用 CSV 字符串格式，字段顺序依赖强，消费端必须严格按列序解析，扩展字段时修改成本高；(3) 无随机种子控制，每次运行结果不可复现。

本项目适配版 Producer（\texttt{realtime\_producer.py}）改为基于种子数据的受控变异：

\begin{verbatim}
source_row = rng.choice(seed_rows)
payload = mutate_row(source_row, event_index, rng)
# JSON 格式输出
handle.write(json.dumps(payload) + "\n")
\end{verbatim}

关键改进包括：以原始 11 字段客户记录为变异基础，保留城市分布、会员等级分布和消费行为之间的业务关联；使用 JSON 键值对格式替代 CSV 字符串，字段以名称标识而非位置，扩展性更强；通过 \texttt{--seed 42} 固定随机种子保证结果可复现；满意度等级由变异后的评分和购买间隔动态推导（而非仅依赖评分），保证字段间逻辑自洽。

\subsubsection{Consumer 端对比}

原始 Consumer（\texttt{kafka\_consumer\_ecommerce\_to\_pg.py}）的处理架构为：

\begin{verbatim}
schema = "customer_id INT, gender STRING, ..."
parsed = kafka_stream_df.select(
    from_csv(col("value"), schema).alias("data")
).select("data.*")
# 年龄分组 → 按 Membership+AgeGroup+City 聚合 → JDBC 写 PostgreSQL
\end{verbatim}

其风险客户识别采用单条件规则：

\begin{verbatim}
days_since_last_purchase > 30
&& satisfaction_level == "Unsatisfied"
\end{verbatim}

输出仅写入 PostgreSQL 数据库，无可视化仪表盘。

本项目适配版 Processor（\texttt{realtime\_processor.py}）在此基础上进行了全面增强：(1) 使用 \texttt{StructType} 显式定义 Schema 替代字符串式 Schema，类型检查更严格；(2) 在原有风险判定条件基础上增加 \texttt{average\_rating < 3.8} 作为替代条件，使灵敏度更高的客户也能被识别；(3) 新增高价值订单标记（\texttt{total\_spend >= 1000}），为会员等级分析和 KPI 统计提供额外维度；(4) 输出从单一数据库写入扩展为 7 类 CSV 文件 + HTML 仪表盘 + Markdown 摘要的完整报告体系；(5) 增加了空批次过滤、批次上限、空闲超时、路径安全检查和 \texttt{try/finally} 资源保护等工程化机制。

\begin{table}[H]
\centering
\caption{原始 Kafka 链路与本项目适配版核心差异}
\begin{tabular}{p{2.8cm}p{4.2cm}p{4.5cm}}
\toprule
对比维度 & 原始 Kafka 链路 & 本项目适配版 \\
\midrule
数据生成 & 完全随机，字段独立生成 & 种子变异，保留业务关联 \\
消息格式 & CSV 字符串（位置依赖） & JSON 键值对（名称标识） \\
Schema 定义 & 字符串式（\texttt{"col INT, ..."}） & \texttt{StructType} 显式字段列表 \\
风险判定 & 单条件（不满意+超30天） & 双条件（增加评分<3.8） \\
高价值订单 & 无 & 有（消费≥1000标记） \\
聚合维度 & 会员×年龄组×城市 & 城市 + 会员类型 + 批次KPI \\
输出形式 & PostgreSQL（JDBC） & CSV + HTML仪表盘 + MD摘要 \\
工程机制 & 基础信号处理 & 空批次过滤、停止策略、路径安全、资源保护 \\
可复现性 & 无随机种子 & \texttt{--seed 42} 固定种子 \\
\bottomrule
\end{tabular}
\end{table}

\subsection{本项目适配与改进}

相较于原始 Kafka 链路，本项目在实时处理模块整体设计中对 Kafka 模式进行了重新适配。改进的目标是使系统既能在教学环境中方便运行，也能够保留面向生产级实时处理架构的扩展能力。

在数据生成方式方面，本项目使用基于种子数据的受控变异方法替代完全随机生成。事件生成器会先读取原始电商客户行为数据集，将其中的客户记录作为种子池，再对消费金额、购买件数、平均评分、距上次购买天数、折扣使用情况和年龄等字段进行合理扰动。通过这种方式生成的交易事件能够保留原始数据集中的统计特征，同时体现实时业务中的动态变化。与完全随机生成相比，种子变异方法能够减少字段之间的矛盾，使客户满意度、评分、购买间隔和消费行为之间保持更合理的关系，从而提高后续实时分析结果的可信度。

在消息格式方面，本项目使用 JSON 格式替代原始实现中较简单的 CSV 字符串格式。CSV 字符串虽然结构紧凑，但字段顺序依赖较强，消费端必须严格按照固定列顺序解析数据。一旦字段顺序发生变化，或者后续需要增加字段，就容易导致解析逻辑出错。JSON 格式以键值对形式表达事件内容，每个字段都具有明确名称，结构更加清晰，也更适合表示复杂事件。Kafka 模式下，事件生成器将每条交易事件序列化为 JSON 消息发送到 Topic；流处理引擎消费消息后，从 Kafka 的 \texttt{value} 字段中取出 JSON 字符串，并使用统一 Schema 解析为结构化 DataFrame。

在处理架构方面，本项目统一了文件流模式和 Kafka 模式。事件生成器可以根据参数选择将事件写入本地 JSON Lines 文件，也可以选择将事件发送到 Kafka Topic；流处理引擎同样可以根据参数选择监听文件目录或消费 Kafka 消息。两种模式虽然输入方式不同，但在进入核心处理逻辑后会被统一转换为相同结构的数据。这样，城市维度聚合、会员维度聚合、风险客户识别、KPI 统计和报表输出等逻辑都可以复用同一套实现，避免为不同输入源重复编写处理代码。

在教学演示方面，本项目将文件流模式作为默认运行方式。Kafka 和 Zookeeper 虽然适合生产级实时数据处理，但对于课程实验而言，会增加环境配置、服务启动和问题排查成本。默认使用文件流模式可以使用户在不启动额外服务的情况下完成完整实时处理链路的演示，包括事件生成、Spark 监听、微批次处理、CSV 输出和仪表盘刷新。这样有利于突出流式处理本身的核心逻辑，也方便在本地环境中复现实验结果。

在扩展能力方面，本项目仍然完整保留 Kafka 模式。Kafka 模式能够支持更接近真实业务系统的事件传输方式，使上游事件生成和下游数据处理之间通过消息队列解耦。当系统未来需要接入真实订单服务、用户行为日志或多源交易事件时，可以将 Kafka Topic 作为统一入口，由 Spark Structured Streaming 进行实时消费和分析。同时，Kafka 模式也为后续扩展实时告警、在线推荐、客户流失预测等模块提供了基础。

本项目的改进可以概括为以下几个方面：
\begin{itemize}
    \item 使用种子数据变异替代完全随机生成，提高模拟事件的数据真实性和业务一致性；
    \item 使用 JSON 格式替代 CSV 字符串格式，使消息字段更清晰，结构扩展更方便；
    \item 统一文件流模式和 Kafka 模式，使不同数据源共用同一套流处理逻辑；
    \item 默认采用文件流模式，降低课程实验和教学演示中的环境复杂度；
    \item 预留 Kafka 模式，使系统具备向生产级消息队列架构扩展的能力。
\end{itemize}

通过这些适配与改进，本项目在原始 Kafka 链路基础上形成了更加完整、清晰和可扩展的实时处理架构。文件流模式保证了实验可运行性，Kafka 模式保证了系统扩展性，两者共同支撑了实时处理模块在不同场景下的使用需求。

\subsection{Docker 编排支持}

为了降低运行环境配置难度，并保证不同组件之间能够稳定协同，本项目通过 Docker Compose 对相关服务进行编排。Docker Compose 可以将 Spark、Kafka、Zookeeper 以及项目代码目录等组件组织到统一的容器化环境中，使系统部署和运行过程更加清晰。对于实时处理模块而言，Docker 编排不仅方便启动基础服务，也有助于统一路径映射、网络配置和依赖环境。

在 Kafka 模式下，系统需要依赖 Kafka Broker 和 Zookeeper 服务。Kafka Broker 负责接收事件生成器发送的消息，并将消息保存到指定 Topic 中；Zookeeper 则用于协调 Kafka 集群元数据和服务状态。在 Docker Compose 配置中，项目为 Kafka 相关服务设置了专门的 profile。默认执行 \texttt{docker-compose up} 时，Kafka 和 Zookeeper 不会自动启动；只有在命令中显式指定 Kafka profile 时，相关服务才会被启用。

这种 profile 设计兼顾了简洁性和扩展性。对于大多数教学演示和本地验证场景，用户只需要运行文件流模式，不需要启动 Kafka 相关容器，因此默认配置保持轻量。对于需要验证 Kafka 模式的场景，用户可以通过启用 Kafka profile 启动 Zookeeper 和 Kafka Broker，从而运行完整的消息队列式实时处理链路。这样的设计避免了在默认运行时引入不必要的服务，也保留了进一步扩展系统的能力。

Kafka Broker 的配置中通常会包含监听地址和对外暴露地址。例如，容器内部服务之间可以使用 \texttt{kafka-broker:29092} 进行通信，而宿主机上的程序可以通过 \texttt{localhost:9092} 访问 Kafka。通过区分容器内部地址和宿主机访问地址，系统能够同时支持容器内 Spark 作业和本地调试程序连接 Kafka。该配置对于 Docker 环境下的 Kafka 使用非常重要，因为 Kafka 消费者和生产者需要根据自身运行位置选择正确的连接地址。

Docker Compose 还通过目录挂载实现宿主机与容器之间的数据共享。例如，项目脚本可以挂载到容器内的 Spark 应用目录，实时处理结果可以从容器内写出到宿主机的 \texttt{Reports/} 目录。这样，用户可以在本地编辑代码、启动容器运行任务，并在本地直接查看生成的 CSV 文件、HTML 仪表盘和 Markdown 摘要。该方式形成了较为完整的开发闭环，提高了实验操作的便利性。

从系统设计角度看，Docker 编排支持使 Kafka 模式不只是代码层面的预留，而是具备实际运行的环境基础。Kafka 模式下，producer 通过 \texttt{--mode kafka} 将事件发送至 Topic，processor 通过 \texttt{--source-type kafka} 从 Topic 消费，两者共用与文件流模式完全一致的数据 Schema 和处理逻辑。通过 \texttt{docker-compose --profile kafka up -d} 即可一键启动包含 Zookeeper、Kafka Broker 和 Spark 集群的完整实时处理环境，无需额外手动配置。用户可以根据需要选择轻量的文件流模式，也可以启用 Kafka profile 运行更接近真实生产架构的模式。默认不启动 Kafka 保证了系统的简洁性，按需启用 Kafka 则保证了系统的扩展性。

综上，Docker Compose 在本项目中承担了运行环境统一和服务按需编排的作用。通过 profile 控制 Kafka 与 Zookeeper 的启用方式，系统既避免了默认运行时的复杂依赖，又为 Kafka 模式提供了完整的容器化支持。这种设计与前文的双模式实时处理架构相配合，使项目在课程演示和工程扩展之间取得了较好的平衡。

\section{实验运行与结果验证}

\subsection{实验环境}

本次实时处理实验的运行环境配置如下：

\begin{itemize}
    \item \textbf{操作系统}：Windows 11 Home China（宿主机），Docker 容器内运行 Ubuntu 22.04 Jammy（基于 \texttt{eclipse-temurin:11-jre-jammy} 镜像）；
    \item \textbf{Docker 版本}：Docker Desktop，使用 Compose 文件格式 v3，编排 Hadoop + Spark + PostgreSQL 集群；
    \item \textbf{Spark 版本}：Apache Spark 3.0.2，基于 Hadoop 3.2 构建，以 Standalone 集群模式运行（1 个 master + 2 个 worker 节点）；
    \item \textbf{Python 版本}：Python 3（容器内通过 APT 安装，配套 numpy、pandas、matplotlib、scipy 等科学计算库）；
    \item \textbf{Hadoop 版本}：Hadoop 3.2.1，包含 NameNode、DataNode、ResourceManager、NodeManager、HistoryServer 五个组件；
    \item \textbf{Kafka 是否启用}：未启用。本次实验使用默认的\textbf{文件流模式}（File Stream Mode），事件生成器将 JSON 事件写入共享卷目录，流处理器通过 Spark Structured Streaming 监听该目录。Kafka 模式（Zookeeper + Kafka Broker）通过 Compose \texttt{profiles: [kafka]} 按需启用，仅在需要演示生产级消息队列链路时启动；
    \item \textbf{输入数据集}：\texttt{E-commerceCustomerBehavior-Sheet1.csv}，经典电商用户行为数据集，包含 Customer ID、Gender、Age、City、Membership Type、Total Spend、Items Purchased、Average Rating、Discount Applied、Days Since Last Purchase、Satisfaction Level 共 11 个字段；
    \item \textbf{输出目录}：容器内路径 \texttt{/opt/spark-results/realtime/}，通过 Docker 卷挂载映射到宿主机 \texttt{./Reports/realtime/}，所有分析结果可直接在宿主机上查看。
\end{itemize}

\subsection{运行流程}

实验运行流程可概括为以下步骤：

\begin{enumerate}
    \item \textbf{启动 Docker 容器环境}：执行 \texttt{docker-compose up -d}，启动 Hadoop 集群组件（NameNode、DataNode、ResourceManager、NodeManager、HistoryServer）、Spark 集群（spark-master、spark-worker-a、spark-worker-b）和 PostgreSQL 数据库（demo-database）。此时不启用 Kafka profile，默认走文件流模式；
    \item \textbf{运行实时事件生成器}：在 spark-master 容器内执行 \texttt{python /opt/spark-apps/realtime\_producer.py --mode files --iterations 5 --batch-size 20 --sleep-seconds 5 --reset-output}，从种子 CSV 数据集中随机抽取记录并施加受控变异，每批生成 20 条事件，间隔 5 秒，共发布 5 批，写入 \texttt{/opt/spark-data/realtime\_input/} 目录；
    \item \textbf{启动流处理引擎}：在 spark-master 容器内执行 \texttt{spark-submit /opt/spark-apps/realtime\_processor.py --source-type files --trigger-seconds 5 --max-batches 5 --idle-timeout-seconds 12}，Spark Structured Streaming 以 5 秒为微批次触发间隔监听输入目录，每个批次触发时执行完整的数据增强、多维聚合、风险识别、KPI 统计和仪表盘渲染流程；
    \item \textbf{等待若干批次事件被处理}：流处理器在后台持续运行，自动检测并处理 producer 陆续写入的 JSON 事件文件。每处理完一个含有数据的批次，控制台打印当前批次编号、事件数量和输出目录路径；
    \item \textbf{查看输出文件}：处理完成后，在宿主机 \texttt{Reports/realtime/} 目录下查看生成的 7 类 CSV 文件（kpi\_history.csv、latest\_batch\_metrics.csv、latest\_events.csv、stream\_events.csv、city\_summary\_latest.csv、membership\_summary\_latest.csv、at\_risk\_customers\_latest.csv）以及 \texttt{dashboard.html} 仪表盘和 \texttt{summary.md} 摘要；
    \item \textbf{对关键指标进行一致性验证}：逐一核对累计事件数、总收入、平均消费、风险客户数量等指标，确认处理结果与输入数据之间满足逻辑一致性约束。
\end{enumerate}

\begin{figure}[H]
\centering
\includegraphics[width=\textwidth]{figures/terminal_processing.png}
\caption{终端运行过程：生产者按批次发布 JSON 事件文件（左），Spark Structured Streaming 流处理引擎监听输入目录并逐批完成处理（右），控制台实时打印批次编号、事件数量和输出路径}
\label{fig:terminal_run}
\end{figure}

\subsection{数据一致性验证}

本次实验运行 1 个微批次，共处理 36 条实时模拟交易事件，完整贯穿了数据解析、增强、聚合、风险识别、KPI 计算和仪表盘渲染的全链路处理流程。系统设计支持多批次连续处理——生产者默认配置为 \texttt{--iterations 5 --batch-size 20}，流处理引擎通过 \texttt{--max-batches 5} 控制批次上限，配合 5 秒微批次触发间隔与空闲超时机制，能够在多个批次间持续进行数据处理与结果更新。本次实验以 36 条事件所形成的单批次结果为基础，从以下方面验证实时处理结果的正确性与一致性：

\begin{itemize}
    \item \textbf{累计事件数}：\texttt{kpi\_history.csv} 中记录的 \texttt{cumulative\_event\_count} 为 36，与 \texttt{stream\_events.csv} 总行数一致；
    \item \textbf{批次收入}：\texttt{kpi\_history.csv} 中 \texttt{total\_revenue} 为 30,575.22，与 \texttt{latest\_events.csv} 中全部 36 条事件的 \texttt{total\_spend} 字段求和结果一致；
    \item \textbf{平均消费}：\texttt{avg\_spend} = 30,575.22 \,/\, 36 = 849.31，与 \texttt{kpi\_history.csv} 记录值完全吻合；
    \item \textbf{折扣订单数}：\texttt{discounted\_orders} = 15，与 \texttt{latest\_events.csv} 中 \texttt{discount\_applied=True} 的计数一致；
    \item \textbf{高价值订单数}：\texttt{high\_value\_orders} = 11，与 \texttt{latest\_events.csv} 中 \texttt{total\_spend} $\geq$ 1000 的事件数一致（11 条全部来自 Gold 会员）；
    \item \textbf{风险客户数}：\texttt{at\_risk\_count} = 9，与 \texttt{at\_risk\_customers\_latest.csv} 的数据行数一致；
    \item \textbf{城市数量}：城市聚合结果包含 6 个城市（Chicago、Houston、Los Angeles、Miami、New York、San Francisco），\texttt{city\_summary\_latest.csv} 共 6 行，与 \texttt{latest\_events.csv} 中 \texttt{city} 列 distinct 值一致；
    \item \textbf{会员类型数量}：会员聚合结果包含 3 种类型（Bronze、Gold、Silver），\texttt{membership\_summary\_latest.csv} 共 3 行，与 \texttt{latest\_events.csv} 中 \texttt{membership\_type} 列 distinct 值一致。
\end{itemize}

\begin{table}[H]
\centering
\caption{数据一致性验证结果}
\begin{tabular}{llll}
\toprule
验证项 & 计算方式 & 实际结果 & 状态 \\
\midrule
累计事件数 & \texttt{stream\_events.csv} 总行数 & 36 & 通过 \\
批次总收入 & \texttt{latest\_events.csv} 中 \texttt{total\_spend} 求和 & 30,575.22 & 通过 \\
平均消费 & 总收入 $/$ 事件数 & 849.31 & 通过 \\
平均评分 & \texttt{latest\_events.csv} 中 \texttt{average\_rating} 均值 & 4.04 & 通过 \\
折扣订单数 & \texttt{discount\_applied=True} 计数 & 15 & 通过 \\
高价值订单数 & \texttt{total\_spend} $\geq$ 1000 计数 & 11 & 通过 \\
风险客户数 & \texttt{at\_risk\_customers\_latest.csv} 行数 & 9 & 通过 \\
城市数 & \texttt{city\_summary\_latest.csv} 行数 & 6 & 通过 \\
会员类型数 & \texttt{membership\_summary\_latest.csv} 行数 & 3 & 通过 \\
\bottomrule
\end{tabular}
\end{table}

此外，还可从各维度聚合结果的业务合理性角度校验：Gold 会员平均消费 1,372.32，显著高于 Silver 的 770.91 和 Bronze 的 479.14；所有 11 笔高价值订单均来自 Gold 会员；San Francisco 和 New York 的平均消费（1,512.31 和 1,204.33）显著高于其他城市。这些结果符合电商业务的直觉预期。

\subsection{流程稳定性验证}

系统在运行过程中针对以下方面进行了稳定性测试：

\begin{itemize}
    \item \textbf{空批次跳过}：当监听目录中无新 JSON 文件到达时，\texttt{process\_batch()} 函数在入口处检测 \texttt{event\_count == 0} 并直接返回，不产生任何空 CSV 文件，也不更新 KPI 计数器。测试确认：在 producer 批次间隔期间，处理器正确跳过了无数据的微批次触发；
    \item \textbf{重置与可复现性}：使用 \texttt{--reset-output} 参数重启处理器时，\texttt{safe\_reset\_output()} 函数首先校验目标目录 basename 为 \texttt{"realtime"}（防止误删其他目录），然后删除整个输出目录并重新生成。固定随机种子（\texttt{--seed 42}）保证了 producer 每次生成相同的事件序列，多次运行结果完全一致；
    \item \textbf{空闲超时自动停止}：配置 \texttt{--idle-timeout-seconds 12} 后，当 producer 停止生成事件超过 12 秒，流处理器的轮询循环检测到 \texttt{datetime.now() - last\_batch\_completed\_at >= timedelta(seconds=12)}，自动调用 \texttt{query.stop()} 终止 Structured Streaming 查询，SparkSession 正常关闭，无僵尸进程残留；
    \item \textbf{追加型 CSV 表头去重}：\texttt{append\_csv()} 函数通过检测目标文件是否存在决定是否写入表头（\texttt{header = not os.path.exists(path)}）。对 \texttt{kpi\_history.csv} 和 \texttt{stream\_events.csv} 两个追加型文件，首次写入时包含列名，后续追加仅写入数据行，不会出现重复表头；
    \item \textbf{仪表盘自动刷新}：每个微批次处理完毕后，\texttt{process\_batch()} 内部调用 \texttt{render\_dashboard(args.output\_dir)}，重新读取所有最新 CSV 文件生成 \texttt{dashboard.html} 和 \texttt{summary.md}。同时，\texttt{main()} 函数的 \texttt{finally} 块中也包含一次 \texttt{render\_dashboard()} 调用，确保即使流处理异常终止，最终状态的仪表盘仍能被保存；
    \item \textbf{异常退出安全}：\texttt{main()} 使用 \texttt{try/finally} 结构包裹整个流查询生命周期，\texttt{finally} 块中依次执行仪表盘渲染和 \texttt{spark.stop()}，保证 SparkSession 资源无论正常还是异常退出都会被释放；
    \item \textbf{路径安全校验}：\texttt{safe\_reset\_output()} 中含 basename 白名单检查，拒绝删除 basename 不为 \texttt{"realtime"} 的目录，防止因参数误配置导致的数据丢失事故。
\end{itemize}

\subsection{系统处理效率}

本实时处理系统的性能设计围绕微批次处理模式展开，各环节的处理效率与数据规模、触发间隔和资源配置相关。

在事件生成端，\texttt{realtime\_producer.py} 的种子变异逻辑为纯内存计算，每次变异操作仅涉及数值浮动、布尔翻转和字段赋值，单条事件生成耗时在毫秒级别。以每批 20 条事件为例，$20 \times (抽取 + 变异 + 序列化)$ 的总耗时远小于批次间隔（默认 5 秒），因此生产者端不会成为链路瓶颈。

在流处理端，\texttt{process\_batch()} 的处理耗时主要取决于批次内事件数量。对于 36 条事件的单批次而言，处理流程涵盖：(1) Spark DataFrame 排序与 Pandas 转换；(2) 城市维度 \texttt{groupBy} 聚合（6 个分组）；(3) 会员维度 \texttt{groupBy} 聚合（3 个分组）；(4) 风险客户过滤与排序；(5) KPI 汇总计算；(6) 7 个 CSV 文件的读写（含 3 个覆盖写入和 2 个追加写入）；(7) HTML 仪表盘和 Markdown 摘要的渲染。在 Docker 容器内单机 Spark Standalone 环境下，该批次可以在秒级时间内完成全部处理并稳定输出结果，远小于微批次触发间隔，为后续处理更多事件预留了充分的性能余量。

在可扩展性方面，本系统通过以下设计保证了处理效率：(1) \texttt{foreachBatch} 机制复用优化后的 DataFrame 操作，避免逐条事件处理带来的调度开销；(2) Spark 的 Catalyst 优化器能够对 \texttt{groupBy}、\texttt{agg}、\texttt{filter} 等操作进行自动查询优化；(3) 快照型 CSV 采用覆盖写入而非全量读取后修改，避免了 I/O 随历史记录增长而恶化；(4) 追加型 CSV 只追加新增行，写入开销与批次事件数成正比而非与累计事件数成正比；(5) HTML 仪表盘为静态文件生成，不依赖 Web 服务器或数据库查询，渲染完成后无额外运行时开销。若后续事件量增大，可通过调整 Spark 的 \texttt{shuffle.partitions}（当前为 4）和 Worker 节点资源进行纵向扩展，也可通过 Kafka 模式下的消费者组实现横向扩展。

\subsection{典型运行结果展示}

为直观展示实时处理链路的输出效果，以下列出本次实验运行产生的关键文件及其内容摘要：

\begin{table}[H]
\centering
\caption{各城市聚合结果（city\_summary\_latest.csv）}
\begin{tabular}{lrrr}
\toprule
城市 & 事件数 & 平均消费（\$） & 购买总件数 \\
\midrule
Chicago & 4 & 507.26 & 39 \\
Houston & 9 & 466.64 & 70 \\
Los Angeles & 7 & 817.38 & 79 \\
Miami & 5 & 705.85 & 69 \\
New York & 5 & 1,204.33 & 77 \\
San Francisco & 6 & 1,512.31 & 135 \\
\bottomrule
\end{tabular}
\end{table}

\begin{table}[H]
\centering
\caption{各会员类型聚合结果（membership\_summary\_latest.csv）}
\begin{tabular}{lrrrr}
\toprule
会员类型 & 事件数 & 平均消费（\$） & 平均评分 & 高价值订单数 \\
\midrule
Bronze & 13 & 479.14 & 3.44 & 0 \\
Gold & 11 & 1,372.32 & 4.68 & 11 \\
Silver & 12 & 770.91 & 4.10 & 0 \\
\bottomrule
\end{tabular}
\end{table}

从以上两表可以观察到：(1) Gold 会员的消费能力和评分均显著高于其他等级，反映了会员等级与实际价值的一致性；(2) San Francisco 和 New York 等一线城市的消费水平明显领先；(3) 全部 11 笔高价值订单集中在 Gold 会员，且其平均评分为 4.68，说明高价值客户同时具有较高的满意度。

\begin{table}[H]
\centering
\caption{风险客户明细（at\_risk\_customers\_latest.csv）}
\begin{tabular}{rrrlrrrl}
\toprule
客户编号 & 城市 & 会员类型 & 消费（\$） & 评分 & 距上次购买（天） & 满意度 \\
\midrule
428 & Miami & Silver & 737.16 & 3.8 & 68 & Unsatisfied \\
380 & Miami & Silver & 716.89 & 4.2 & 64 & Unsatisfied \\
428 & Miami & Silver & 646.01 & 3.9 & 59 & Unsatisfied \\
141 & Miami & Silver & 807.12 & 3.9 & 54 & Unsatisfied \\
213 & Miami & Silver & 622.05 & 4.3 & 51 & Unsatisfied \\
295 & Chicago & Bronze & 452.64 & 3.6 & 43 & Unsatisfied \\
133 & Chicago & Bronze & 518.90 & 4.0 & 39 & Unsatisfied \\
393 & Houston & Bronze & 493.33 & 3.7 & 34 & Neutral \\
235 & Chicago & Bronze & 613.78 & 3.7 & 33 & Neutral \\
\bottomrule
\end{tabular}
\end{table}

从风险客户明细表可以观察到：风险最高的客户（距上次购买 68 天）集中在 Miami，且全部 5 名 Miami 风险客户的距上次购买天数均在 50 天以上、满意度均为 Unsatisfied，属于最需关注的流失高风险群体。值得注意的是，客户编号 428 在风险列表中出现了两次（68 天和 59 天），这是因为该客户的不同交易事件分别触发了风险判定条件，反映了同一客户可能在多个时间点持续处于风险状态。Houston 的客户 393 虽然距上次购买仅 34 天且满意度为 Neutral，但由于其评分仅为 3.7（低于 3.8 阈值），按照增强后的风险判定规则也被正确识别——这体现了实时处理在批处理单条件规则基础上引入"评分低于 3.8"作为替代条件所带来的检测灵敏度提升。

KPI 摘要卡片显示：本批次处理 36 条事件，总收入 \$30,575.22，平均单笔消费 \$849.31，风险客户 9 名，平均评分 4.04。这些指标全部通过 HTML 仪表盘（\texttt{dashboard.html}）以卡片、柱状图和表格三种可视化形式实时呈现。

\begin{figure}[H]
\centering
\includegraphics[width=\textwidth]{figures/dashboard_results1.png}
\caption{KPI 摘要卡片、城市平均消费柱状图与会员事件分布}
\label{fig:dashboard_results1}
\end{figure}

\begin{figure}[H]
\centering
\includegraphics[width=\textwidth]{figures/dashboard_results2.png}
\caption{批次 KPI 历史表与风险客户 Top 列表}
\label{fig:dashboard_results2}
\end{figure}

\begin{figure}[H]
\centering
\includegraphics[width=\textwidth]{figures/dashboard_results3.png}
\caption{最新批次事件明细表}
\label{fig:dashboard_results3}
\end{figure}

\begin{figure}[H]
\centering
\includegraphics[width=\textwidth]{figures/dashboard_results4.png}
\caption{会员快照与完整数据表格}
\label{fig:dashboard_results4}
\end{figure}


\section{与批处理链路的对比分析}

\subsection{批处理与实时处理的设计差异}

本项目的批处理链路以 \texttt{apps/batch\_analysis.py} 为核心，实时处理链路以 \texttt{apps/realtime\_processor.py} 为核心，二者虽然共用同一份原始数据集，但在架构设计、数据流转和分析目标上存在根本差异。下表从六个关键维度对两条链路进对比：

\begin{longtable}{p{2.4cm}p{5.2cm}p{5.2cm}}
\caption{批处理分析与实时处理分析对比} \\
\toprule
维度 & 批处理分析 & 实时处理分析 \\
\midrule
\endfirsthead
\caption{批处理分析与实时处理分析对比（续）} \\
\toprule
维度 & 批处理分析 & 实时处理分析 \\
\midrule
\endhead
\bottomrule
\endfoot
数据范围 & 全量历史数据（读取 HDFS 上的完整 CSV 文件，经 \texttt{na.drop()} 清洗） & 最新交易事件流（producer 从种子数据变异生成，每批 20 条） \\
\midrule
处理方式 & 一次性全量计算：依次执行客户分群、K-Means、逻辑回归三个独立阶段，每个阶段遍历全量数据 & 微批次增量计算：每个批次独立触发完整流程（增强 → 聚合 → 风险识别 → KPI → 仪表盘），批次间通过 state 字典跟踪累计状态 \\
\midrule
分析重点 & \begin{itemize}
    \item 多维度客户分群（Membership Type × Age Group × City）
    \item K-Means 聚类（4 类，Silhouette 评估）
    \item 满意度预测（Logistic Regression，80/20 划分）
    \item 城市级消费汇总
\end{itemize} & \begin{itemize}
    \item 实时 KPI 看板（总收入、平均消费、订单构成）
    \item 城市级实时聚合
    \item 会员类型实时聚合
    \item 风险客户动态识别（双重判定条件）
\end{itemize} \\
\midrule
时效性 & 静态离线，运行一次产出最终报告，适合定期离线分析 & 动态近实时，5 秒微批次间隔持续产出，适合运营监控场景 \\
\midrule
技术栈 & \begin{itemize}
    \item PySpark DataFrame API
    \item MLlib（KMeans、Logistic Regression、Pipeline）
    \item StringIndexer、VectorAssembler
    \item MulticlassClassificationEvaluator
\end{itemize} & \begin{itemize}
    \item Spark Structured Streaming
    \item \texttt{readStream.json()} / \texttt{readStream.format("kafka")}
    \item \texttt{foreachBatch} 回调
    \item 原生 Pandas + CSV I/O
\end{itemize} \\
\midrule
输出形式 & 7 个 CSV 文件 + 1 个 Markdown 摘要 + 静态 HTML 分析仪表盘 + 4 张 PNG 图表（由 \texttt{render\_analysis\_dashboard.py} 渲染） & 7 类 CSV 文件（含追加型时间序列）+ 1 个 Markdown 摘要 + 实时刷新的 HTML 仪表盘（含内联 CSS 条形图和卡片） \\
\midrule
数据源 & HDFS \texttt{hdfs://namenode:9000/E-commerceCustomerBehavior-Sheet1.csv} & 文件目录 \texttt{/opt/spark-data/realtime\_input/} 或 Kafka Topic \texttt{ecommerce-transactions} \\
\midrule
启动方式 & \texttt{spark-submit batch\_analysis.py}（一次性执行完毕后退出） & \texttt{spark-submit realtime\_processor.py}（长期驻留，达到 max-batches 或空闲超时后自动退出） \\
\bottomrule
\end{longtable}

\subsection{风险客户识别逻辑对比}

两条链路均包含风险客户识别功能，但判定逻辑的精细度存在差异：

\begin{itemize}
    \item \textbf{批处理}（\texttt{apps/batch\_analysis.py:53-56}）：采用单一判定条件——
    \begin{verbatim}
    at_risk = df.filter(
        (col("Days Since Last Purchase") > 30)
        & (col("Satisfaction Level") == "Unsatisfied")
    )
    \end{verbatim}
    即同时满足"距上次购买超过 30 天"和"满意度为 Unsatisfied"两个条件。

    \item \textbf{实时处理}（\texttt{apps/realtime\_processor.py:148-160}）：在批处理条件基础上进行了增强——
    \begin{verbatim}
    is_at_risk = when(
        (col("days_since_last_purchase") > 30)
        & (
            (col("satisfaction_level") == lit("Unsatisfied"))
            | (col("average_rating") < lit(3.8))
        ),
        lit(True),
    ).otherwise(lit(False))
    \end{verbatim}
    增加了"评分低于 3.8"作为替代条件，即：距上次购买超过 30 天 且（不满意 或 评分低于 3.8）= 风险客户。这意味着即使满意度字段标注为 Neutral，只要评分偏低且长期未购买，也会被识别为风险客户，覆盖面更广、灵敏度更高。
\end{itemize}

\subsection{两条链路的互补关系}

批处理链路能够回答\textbf{"过去发生了什么"}，适合分析整体客户特征、长期消费趋势、用户分群结构和满意度分类规律；实时处理链路能够回答\textbf{"现在正在发生什么"}，适合监控最新交易变化、识别当前风险客户和跟踪实时业务指标。两者的具体互补关系体现在：

\begin{enumerate}
    \item \textbf{分析层次的互补}：批处理提供宏观画像（如 K-Means 聚类揭示的 4 类客户群体、Logistic Regression 模型对满意度的可预测性结论），实时处理提供微观监测（如当前批次某城市突然出现的多笔高价值订单、某客户满意度骤降触发风险告警）；
    \item \textbf{时间粒度的互补}：批处理覆盖整个数据集的全部历史周期，适合生成周报、月报级别的总结性结论；实时处理以 5 秒为粒度更新，适合运营人员需要即时掌握的业务波动；
    \item \textbf{数据流转的互补}：批处理读取 HDFS 上的静态 CSV 文件，实时处理监听文件目录的动态变化——两者共享同一套 Docker 卷挂载体系（\texttt{./data → /opt/spark-data}，\texttt{./Reports → /opt/spark-results}），但分别对应"存量数据"和"增量数据"的分析需求；
    \item \textbf{输出形态的互补}：批处理产生可嵌入报告的高分辨率 PNG 图表（figure1.png $\sim$ figure4.png）和模型评估指标（Silhouette 系数、分类准确率），实时处理产生可嵌入运营大屏的实时 HTML 仪表盘（\texttt{dashboard.html}）和时序化 KPI 历史（\texttt{kpi\_history.csv}）。
\end{enumerate}

\subsection{工程实现的共性}

尽管两条链路在分析模式上截然不同，但工程层面共享以下设计原则：

\begin{itemize}
    \item \textbf{统一的数据 Schema}：批处理的 CSV 读取与实时处理的 JSON Schema 定义（\texttt{build\_schema()} 函数，共 13 个字段）都覆盖同一份数据集的核心字段，保证了分析口径的一致性；
    \item \textbf{统一的输出目录结构}：批处理输出至 \texttt{Reports/batch\_analysis/}，实时处理输出至 \texttt{Reports/realtime/}，均由 Docker 卷挂载映射到宿主机，都包含 CSV 数据文件 + Markdown 摘要 + HTML 仪表盘的完整组合；
    \item \textbf{统一的容器化执行环境}：两条链路均在 spark-master 容器内以 \texttt{spark-submit} 方式运行，共享同一个 Spark 集群，无需额外的环境切换；
    \item \textbf{防御性编程}：批处理的 \texttt{na.drop()} 与实时处理的空批次跳过、\texttt{safe\_reset\_output()} 路径校验，均体现了对异常数据的防御性处理。
\end{itemize}

综上，两条链路共同构成了"历史洞察 + 实时监测"的完整分析体系，使本项目不仅具备对电商用户行为的深度挖掘能力，也具备对业务动态变化的实时感知能力。


\section{关键技术总结}

本项目的实时处理与流式分析模块围绕电商客户行为数据展开，重点解决了实时事件生成、流式数据接入、微批次处理、结果输出和系统稳定性等问题。与单纯的离线批处理分析相比，该模块更强调数据持续到达、结果动态更新和业务状态实时监测。通过事件生成器、流处理引擎和报表生成器之间的协同，系统形成了一条较为完整的实时数据分析链路。本节从事件生成、处理架构、微批次分析、结果写入、可视化展示以及稳定性设计等方面，对本模块中的关键技术进行总结。

\subsection{基于种子变异的事件生成}

事件生成是实时处理链路的输入基础。为了模拟电商平台中持续产生的交易事件，本项目没有采用完全随机生成数据的方式，而是基于原始电商客户行为数据进行种子变异。系统先从原始 CSV 数据集中读取客户记录，并将这些记录作为事件生成的种子池。每次生成新事件时，程序从种子池中抽取一条原始记录，再对消费金额、购买件数、平均评分、距上次购买天数、折扣使用情况和年龄等字段进行受控扰动。

这种设计的优势在于能够在真实性和动态性之间取得平衡。完全随机生成虽然实现简单，但容易产生不符合业务逻辑的数据，例如评分与满意度不匹配、消费金额过低或过高、客户属性与交易行为之间缺乏关联等。基于种子数据变异则可以保留原始数据集中的城市分布、会员等级分布和消费行为特征，使生成事件更接近真实电商业务场景。

同时，受控变异机制能够为静态历史数据引入动态变化。原始数据集本身是固定的，如果直接重复使用原始记录，无法体现实时数据流中事件不断变化的特点。通过在合理范围内调整关键字段，系统既保留了原始数据的统计特征，又使每个批次的交易事件具有一定差异，从而能够支撑后续实时 KPI 统计、风险客户识别和会员行为分析。

此外，事件生成器还为每条交易事件添加唯一事件编号和事件时间，并根据变异后的平均评分与距上次购买天数重新推导满意度等级。该设计保证了事件字段的完整性和逻辑自洽性，也使后续流处理模块可以按照事件时间和事件编号进行排序、归档和追踪。固定随机种子的使用进一步提高了实验结果的可复现性，便于课程项目中的多次运行、结果验证和报告撰写。

\subsection{统一的双模式实时处理架构}

本项目在实时处理架构中同时支持文件流模式和 Kafka 模式。文件流模式主要面向教学演示和本地实验场景，事件生成器将每个批次的交易事件写入 JSON Lines 文件，Spark Structured Streaming 监听输入目录并自动读取新增文件。该模式不依赖额外消息队列服务，运行环境简单，便于调试和展示。

Kafka 模式则面向更接近生产环境的实时数据处理场景。事件生成器可以将交易事件序列化为 JSON 消息并发送到 Kafka Topic，流处理引擎从 Kafka 中消费消息，再将消息内容解析为结构化事件。Kafka 作为消息队列能够实现事件生产端与消费端之间的解耦，适合多生产者、多消费者以及分布式实时处理场景。

两种模式虽然在输入方式上不同，但系统通过统一 Schema 和统一处理逻辑实现了架构上的一致性。文件流模式下，系统直接读取 JSON 文件并应用 Schema；Kafka 模式下，系统先从 Kafka 消息的 \texttt{value} 字段中提取 JSON 字符串，再使用同一套 Schema 解析。完成解析后，两种模式都会进入相同的数据增强、聚合统计、风险客户识别和结果输出流程。

这种双模式架构降低了代码重复度，也提高了系统的可扩展性。默认文件流模式使系统能够在课程环境中快速运行，预留 Kafka 模式则为后续扩展到生产级消息队列架构提供基础。通过参数切换运行模式，系统能够在不同环境下灵活适配，体现了较好的工程设计思想。

\subsection{基于 foreachBatch 的微批次分析}

本项目使用 Spark Structured Streaming 的 \texttt{foreachBatch} 机制实现微批次分析。Structured Streaming 会将持续到达的数据流按照触发间隔划分为多个微批次，每个微批次都可以被视为一个静态 DataFrame。通过 \texttt{foreachBatch}，系统可以在每个批次到达时调用自定义处理函数，对当前批次执行完整的数据分析流程。

在本项目中，每个非空微批次都会依次完成事件留存、城市维度聚合、会员维度聚合、风险客户识别、KPI 指标汇总和仪表盘刷新等步骤。事件留存用于保存当前批次明细和累计事件历史；城市维度聚合用于统计各城市的事件数量、平均消费和购买总件数；会员维度聚合用于统计不同会员等级的事件数量、平均消费、平均评分和高价值订单数量；风险客户识别用于筛选可能存在流失风险的客户；KPI 汇总用于记录当前批次的总体业务表现。

\texttt{foreachBatch} 的优势在于能够在流式处理框架中复用批处理式 DataFrame 操作。系统可以使用 \texttt{groupBy}、\texttt{agg}、\texttt{filter}、\texttt{orderBy} 和 \texttt{select} 等常规操作完成复杂分析，而不需要为每条事件单独编写处理逻辑。这种方式既保留了 Spark DataFrame API 的表达能力，也使代码结构更加清晰。

微批次分析机制还便于控制实验流程。系统可以设置触发间隔、批次上限和空闲超时，使实时处理任务在完成指定批次数或长时间没有新事件时自动结束。对于课程项目而言，这种机制能够在较短时间内完整展示实时事件生成、流处理计算和结果输出过程，同时避免流处理任务无限运行。

\subsection{快照与历史分离的写入策略}

实时处理模块在输出设计上区分了快照型文件和历史型文件。快照型文件表示系统当前最新状态，通常采用覆盖写入方式；历史型文件用于记录处理过程中的完整变化，通常采用追加写入方式。这种写入策略使系统既能够快速展示最新结果，也能够保留完整的批次级处理历史。

快照型文件包括 \texttt{latest\_batch\_metrics.csv}、\texttt{latest\_events.csv}、\texttt{city\_summary\_latest.csv}、\texttt{membership\_summary\_latest.csv} 和 \texttt{at\_risk\_customers\_latest.csv} 等。这些文件每个批次都会被最新结果覆盖，因此始终反映最近一次处理后的状态。HTML 仪表盘读取这些快照文件时，不需要从大量历史记录中筛选最新数据，从而简化了展示逻辑，也提高了页面生成效率。

历史型文件包括 \texttt{kpi\_history.csv} 和 \texttt{stream\_events.csv} 等。\texttt{kpi\_history.csv} 每处理一个批次就追加一条 KPI 记录，可以形成批次级时间序列；\texttt{stream\_events.csv} 则保存所有已处理事件，用于记录完整实时事件历史。通过追加写入，系统能够保留从任务开始到结束的全部处理轨迹，便于后续复盘和验证。

除快照和历史文件外，系统还会将每个批次的事件单独保存到 \texttt{events/batch\_XXX.csv} 文件中。该设计使每个批次都有独立归档，便于对单个批次进行追踪。当某个批次的统计结果出现异常时，可以直接查看对应批次的事件明细，并与聚合结果进行比对。

快照与历史分离的写入策略体现了实时分析系统中的分层输出思想。快照文件服务于实时展示，历史文件服务于过程追踪，批次归档服务于调试验证。三者结合后，系统既能满足仪表盘快速读取最新结果的需求，也能保留可追溯的数据处理过程。

\subsection{无外部依赖的 HTML 仪表盘}

为了增强实时处理结果的展示效果，本项目设计了 HTML 仪表盘。该仪表盘由报表生成器读取流处理输出的 CSV 文件后自动生成，主要展示处理批次数、累计事件数、最新批次收入、风险客户数量、城市平均消费、会员类型分布、KPI 历史、风险客户列表、最新事件明细和会员快照等内容。

仪表盘采用静态 HTML 文件形式，不依赖前端框架、Web 服务器或外部可视化库。样式部分通过内联 CSS 实现，数据展示部分使用原生 HTML 表格和简单柱状图结构完成。用户只需要在浏览器中打开 \texttt{dashboard.html}，即可查看实时分析结果。这种设计降低了系统运行和展示的环境要求，非常适合课程实验和本地演示。

无外部依赖的设计也提高了结果文件的可移植性。由于仪表盘本身是一个独立 HTML 文件，用户可以直接复制、提交或嵌入实验材料中，不需要额外打包 JavaScript、CSS 或图片资源。对于期末报告和项目展示而言，这种形式更容易管理，也更便于检查。

虽然该仪表盘不是复杂的交互式前端系统，但它已经能够覆盖实时处理模块的核心展示需求。摘要卡片用于快速查看系统状态，城市和会员图表用于展示聚合结果，KPI 历史表用于观察批次变化，风险客户列表用于突出业务重点，最新事件明细表用于追踪输入数据。通过这些内容，仪表盘能够较完整地呈现实时处理链路的运行结果。

\subsection{安全与稳定性设计}

实时处理模块在实现过程中加入了多项安全与稳定性机制，以保证系统能够在实验环境中可靠运行。由于实时处理任务涉及文件写入、目录清空、批次处理和长时间监听等操作，如果缺少必要保护，容易出现空文件、重复写入、误删目录或任务无法结束等问题。

在路径安全方面，系统在清空输出目录时加入了目录名检查。只有当目标目录明确为实时处理输出目录时，程序才会执行删除操作。如果用户误传了项目根目录或其他重要路径，系统会拒绝删除并抛出异常。该机制能够有效避免因参数错误导致的文件误删，提高系统运行安全性。

在批次处理稳定性方面，系统设计了空批次过滤机制。当 Spark Structured Streaming 触发批次但没有读取到新事件时，系统会直接跳过该批次，不执行后续写文件和报表刷新操作。这样可以避免生成无意义的空结果文件，也可以减少对仪表盘展示的干扰。

在写入稳定性方面，系统区分覆盖写入和追加写入，并在追加 CSV 文件时处理表头问题。历史型文件只在首次写入时添加表头，后续追加记录时不重复写表头，从而保证 CSV 文件结构正确。快照型文件则每次覆盖写入，保证内容始终代表最新状态。

在运行控制方面，系统设计了批次上限和空闲超时两类停止条件。批次上限用于控制实验处理规模，当系统处理到指定数量批次后自动结束；空闲超时用于处理事件生成器已经停止但流处理引擎仍在等待的情况，当长时间没有新事件到达时自动关闭任务。这两类机制使系统适合课程项目中的有限时长演示，避免流处理任务无限运行。

综合来看，安全与稳定性设计提高了实时处理模块的健壮性。路径检查防止误删，空批次过滤防止无效输出，写入策略保证文件结构稳定，停止策略保证任务可控结束。这些机制虽然不直接改变分析结果，但对于一个能够反复运行、便于展示和可维护的数据处理系统而言非常重要。

\section{系统特点与后续扩展方向}

\subsection{当前系统的设计定位}

本项目实现的实时处理与流式分析模块定位为课程实验与教学演示场景下的完整实时分析系统。在架构设计上，系统优先考虑了教学环境中的易用性、可复现性和展示效果。文件流模式作为默认运行方式，降低了环境依赖；事件生成的种子变异策略保证了模拟数据的业务合理性；无外部依赖的 HTML 仪表盘使结果展示开箱即用；快照与历史分离的输出策略兼顾了实时展示与过程追踪。通过上述设计选择，系统能够在有限的课程实验时间内完成从事件生成到可视化展示的完整演示流程。

系统的部分设计选择体现了教学场景与真实生产环境之间的自然差异，这些差异同时也为系统的进一步扩展提供了明确的方向。

\subsection{扩展方向}

基于当前系统已有的架构基础，后续可以从以下维度进行扩展，使系统逐步向更完整的实时数据分析平台演进。

在流式计算能力方面，可以在当前微批次处理基础上引入滑动窗口和滚动窗口分析。当前系统以微批次为基本统计单位，每个批次独立完成聚合计算。若引入 Spark Structured Streaming 的窗口操作，系统可以统计固定时间范围或滑动时间范围内的交易指标，例如最近 5 分钟总收入、最近 10 分钟风险客户数量、最近 1 小时不同城市消费趋势等。窗口分析能够提供更丰富的时间维度视角，与批次级快照形成互补。

在客户风险识别方面，可以在当前规则引擎基础上接入机器学习模型。当前基于人工阈值的风险判定规则具有可解释性强、部署简单和实时响应快等优点。在此基础上，后续可以引入客户流失预测模型，将购买频率、消费金额变化趋势、评分变化轨迹、会员等级、折扣使用习惯等多维特征纳入模型输入，输出客户流失概率或风险分数。将规则引擎与机器学习模型结合，可以实现"规则快速筛选 + 模型精细评分"的分层风险识别体系。

在结果存储方面，可以在当前 CSV 文件输出基础上增加数据库写入支持。当前 CSV 文件适合本地查看和实验验证，便于课程报告中的数据提取和使用。当系统需要支持更长期的数据存储和更复杂的查询时，可以将批次 KPI、事件明细、城市聚合结果和风险客户记录同步写入 PostgreSQL 或 ClickHouse 等数据库。本项目的 Docker Compose 已包含 PostgreSQL 服务（demo-database），原始 Kafka Consumer 中也已有 JDBC 写入的参考实现，因此该扩展具有现成的环境和技术基础。

在仪表盘展示方面，可以在当前静态 HTML 仪表盘基础上增强交互能力。当前仪表盘采用内联 CSS 和原生 HTML 实现，无需任何外部依赖，适合快速查看和结果提交。在此基础上，后续可以引入 ECharts、D3.js 等前端可视化库，增加按城市筛选、会员类型过滤、时间范围切换、数据排序和分页等交互功能，使仪表盘从"查看型"升级为"探索型"分析工具。

在 Kafka 模式部署方面，当前系统已完成完整的代码适配和 Docker 编排。Producer 和 Processor 均通过参数切换支持 Kafka 模式，Docker Compose 通过 profile 机制按需启用 Zookeeper 和 Kafka Broker。后续可以进一步完善 Topic 分区策略、消费者组配置、消息偏移管理和容错重试机制，使其更接近生产级消息队列架构。Checkpoint 机制已在代码中配置完毕，能够支持流处理任务的重启恢复。

在异常检测与告警方面，可以在当前的可视化监控基础上增加自动提醒机制。当前仪表盘已经展示了风险客户数量、最新收入、平均评分等核心指标，用户可以通过查看仪表盘了解系统状态。后续可以为关键指标设置动态阈值或异常检测规则，当批次风险客户数量超过设定比例、平均评分显著下降或交易收入异常波动时，通过日志标记、仪表盘高亮或消息通知等方式提醒用户。

在系统工程能力方面，可以在当前已有的安全机制基础上进一步增强。当前系统已具备路径安全校验、空批次过滤、批次上限、空闲超时和 \texttt{try/finally} 资源保护等机制。后续可以增加结构化日志记录（包含批次处理耗时、事件数量、输出路径和异常信息），支持 YAML/JSON 配置文件管理运行参数，以及提供更细粒度的错误分类与恢复策略。

通过以上扩展，实时处理模块可以在保持教学可用性的同时，逐步发展为功能更完整的实时数据分析平台。窗口统计可增强流式分析的表达能力，机器学习可提升风险识别的智能性，数据库存储可支撑长期查询，交互式仪表盘可改善用户体验，Kafka 生产级部署可增强系统扩展能力，异常告警可提升实时监控的实用价值。这些方向并非对当前系统不足的修补，而是在现有坚实基础上向更高目标的自然延伸。

\section{结论}

本文围绕电商客户行为分析系统中的实时处理与流式分析模块展开设计与实现，构建了一条从模拟交易事件生成、流式数据接入、微批次处理到结果展示输出的完整数据链路。该模块以原始电商客户行为数据为基础，通过受控变异方式生成新的实时交易事件，使事件数据既保留原始数据集的统计特征，又能够体现实时业务中的动态变化。相比完全随机生成方式，基于种子数据的事件生成方法提高了模拟数据的真实性和业务解释性。

在流处理实现方面，本文采用 Spark Structured Streaming 作为核心计算框架，将持续到达的交易事件抽象为不断增长的数据流，并通过微批次机制进行周期性处理。系统同时支持文件流模式和 Kafka 模式。文件流模式降低了课程实验和本地演示的环境复杂度，便于快速运行和验证；Kafka 模式则为系统扩展到生产级消息队列架构提供了基础。通过统一 Schema 设计，两种模式在进入流处理阶段后能够共用同一套数据增强、聚合统计和结果输出逻辑。

在业务分析方面，系统围绕电商客户行为场景设计了多类实时分析指标。流处理引擎在原始事件字段基础上新增风险客户标记和高价值订单标记，并在每个微批次中完成城市维度聚合、会员维度聚合、风险客户识别和 KPI 指标汇总。城市维度分析可以反映不同地区的最新消费表现，会员维度分析可以比较不同会员等级的消费能力和评分情况，风险客户识别则可以帮助发现可能存在流失风险或需要运营关注的客户。

在结果输出方面，系统生成了多种形式的分析结果，包括 CSV 结构化文件、Markdown 摘要和 HTML 实时仪表盘。CSV 文件用于保存批次事件明细、全量事件历史、KPI 时间序列、城市聚合结果、会员聚合结果和风险客户明细，保证处理结果可追踪、可复核。Markdown 摘要用于快速查看当前实时处理状态，HTML 仪表盘则通过核心指标卡片、图表和表格直观展示实时分析结果。通过快照型文件覆盖写入和历史型文件追加写入相结合的策略，系统同时满足了最新结果展示和历史过程追踪的需求。

实验验证表明，该实时处理链路能够稳定处理模拟交易事件，并正确生成各类实时分析结果。系统能够在事件生成器持续写入数据后，由 Spark Structured Streaming 自动监听新增事件并触发微批次处理，最终刷新 CSV 文件、Markdown 摘要和 HTML 仪表盘。与批处理模块相比，实时处理模块更加关注最新业务动态、批次级 KPI 变化和风险客户监测；批处理模块则更适合进行全量历史统计和长期行为分析。两者结合后，电商客户行为分析系统既能够回答“过去发生了什么”，也能够回答“当前正在发生什么”。

总体而言，本项目验证了 Spark Structured Streaming 在电商实时分析场景中的可行性，也展示了文件流模式、Kafka 模式、微批次处理、风险客户识别和实时仪表盘生成等技术在课程项目中的综合应用价值。整体架构已经为后续扩展提供了良好基础。未来可以进一步引入时间窗口聚合、机器学习预测、数据库存储、交互式可视化和异常告警机制，使系统逐步从实验性实时处理模块发展为更加完整的电商实时数据分析平台。



\end{document}